\chapter{Integral Riemann} \label{int:chapter}

%%%%%%%%%%%%%%%%%%%%%%%%%%%%%%%%%%%%%%%%%%%%%%%%%%%%%%%%%%%%%%%%%%%%%%%%%%%%%%

\section{Integral Riemann}
\label{sec:rint}

%mbxINTROSUBSECTION

\sectionnotes{1,5 kuliah}

Integral merupakan cara untuk \myquote{menjumlahkan} nilai-nilai suatu fungsi.
Mahasiswa kalkulus terkadang mencampuradukkan
\emph{integral} dan \emph{antiturunan}.
Secara informal, integral tidak lain adalah luas di bawah kurva.
Bahwa kita dapat menghitung antiturunan menggunakan integral merupakan hasil
tak sepele yang harus kita buktikan.
Kita akan mendefinisikan \emph{integral Riemann}%
\footnote{Dinamai menurut matematikawan Jerman
\href{https://en.wikipedia.org/wiki/Riemann}{Georg Friedrich Bernhard Riemann}
(1826--1866).}
menggunakan integral Darboux%
\footnote{Dinamai menurut matematikawan Prancis
\href{https://en.wikipedia.org/wiki/Darboux}{Jean-Gaston Darboux} (1842--1917).},
suatu definisi yang ekuivalen tetapi secara teknis lebih sederhana.

\subsection{Partisi serta integral bawah dan atas}

Kita ingin mengintegralkan suatu fungsi terbatas yang didefinisikan pada interval $[a,b]$.
Terlebih dahulu kita mendefinisikan dua integral bantu yang terdefinisi untuk semua
fungsi terbatas.  Barulah setelah itu kita dapat membicarakan integral Riemann dan
fungsi-fungsi yang dapat diintegralkan secara Riemann.

\begin{defn}
Suatu \emph{\myindex{partisi}} $P$ dari $[a,b]$ adalah
himpunan berhingga yang terdiri atas bilangan-bilangan $\{ x_0,x_1,x_2,\ldots,x_n \}$ sedemikian sehingga
\begin{equation*}
a = x_0 < x_1 < x_2 < \cdots < x_{n-1} < x_n = b .
\end{equation*}
Kita tulis
\begin{equation*}
\Delta x_i \coloneqq x_i - x_{i-1} .
\end{equation*}
Misalkan $f \colon [a,b] \to \R$ merupakan fungsi terbatas dan $P$ merupakan partisi dari
$[a,b]$.
Definisikan
\begin{align*}
& m_i \coloneqq \inf \, \bigl\{ f(x) : x_{i-1} \leq x \leq x_i \bigr\} , &
& M_i \coloneqq \sup \, \bigl\{ f(x) : x_{i-1} \leq x \leq x_i \bigr\} , \\
& L(P,f) \coloneqq
\sum_{i=1}^n m_i \Delta x_i , &
& U(P,f) \coloneqq
\sum_{i=1}^n M_i \Delta x_i .
\avoidbreak
\end{align*}
Kita menyebut $L(P,f)$ sebagai \emph{\myindex{jumlah Darboux bawah}} dan
$U(P,f)$ sebagai \emph{\myindex{jumlah Darboux atas}}\index{jumlah Darboux}.
\glsadd{not:lowerdarbouxsum}
\glsadd{not:upperdarbouxsum}
\end{defn}

Gagasan geometris jumlah Darboux diilustrasikan dalam
\figureref{darbouxfig}.  Jumlah bawah merupakan luas persegi panjang
berarsir, sedangkan jumlah atas merupakan luas keseluruhan persegi panjang,
yakni bagian berarsir dan tidak berarsir.  Lebar persegi panjang ke-$i$ adalah $\Delta x_i$,
tinggi persegi panjang berarsir adalah $m_i$, dan tinggi
keseluruhan persegi panjang adalah $M_i$.

\begin{myfigureht}
\myincludegraphics{darbouxfig}{%
Diagram grafik suatu fungsi dengan garis hitam tebal, seluruhnya berada di atas
sumbu x.
Domain dibagi menjadi 8 interval dan
titik-titik ujungnya ditandai, dari kiri ke kanan, x subskrip 0 hingga x subskrip 8.
Interval antara x subskrip 4 dan x subskrip 5 ditandai memiliki lebar
Delta x subskrip 5.
Di atas setiap interval terdapat dua persegi panjang.  Sebagai contoh,
terdapat persegi panjang berarsir dengan lebar Delta x subskrip 5
yang sisi bawahnya tepat berupa interval dari
x subskrip 4 hingga x subskrip 5 dan tingginya m subskrip 5, yaitu nilai minimum
fungsi pada interval tersebut.  Persegi panjang itu diperpanjang ke atas tanpa arsiran hingga
M subskrip 5, yaitu nilai maksimum fungsi pada interval tersebut.}
\caption{Contoh jumlah Darboux.\label{darbouxfig}}
\end{myfigureht}

\begin{prop} \label{sumulbound:prop}
Misalkan $f \colon [a,b] \to \R$ merupakan fungsi terbatas.  Misalkan $m, M \in \R$ 
sedemikian sehingga untuk semua $x \in [a,b]$ berlaku $m \leq f(x) \leq M$.  Maka, untuk setiap partisi $P$
dari $[a,b]$,
\begin{equation}
\label{sumulbound:eq}
m(b-a) \leq
L(P,f) \leq U(P,f)
\leq M(b-a) .
\end{equation}
\end{prop}

\begin{proof}
Misalkan $P$ merupakan partisi dari $[a,b]$.  Perhatikan bahwa untuk semua $i$,
berlaku $m \leq m_i \leq M_i \leq M$.  Kita juga memiliki
$\sum_{i=1}^n \Delta x_i = (b-a)$.  Oleh karena itu,
\begin{multline*}
m(b-a) =
m \left( \sum_{i=1}^n \Delta x_i \right)
=
\sum_{i=1}^n m \Delta x_i
\leq
\sum_{i=1}^n m_i \Delta x_i 
\\
\leq
\sum_{i=1}^n M_i \Delta x_i
\leq
\sum_{i=1}^n M \Delta x_i 
=
M \left( \sum_{i=1}^n \Delta x_i \right)
=
M(b-a) .
\end{multline*}
Jadi, kita memperoleh \eqref{sumulbound:eq}.  Secara khusus, himpunan semua jumlah
bawah dan himpunan semua jumlah atas bersifat terbatas.
\end{proof}

\begin{defn}
Karena himpunan semua jumlah Darboux bawah maupun atas terbatas, kita definisikan
\begin{align*}
& \underline{\int_a^b} f(x)\,dx \coloneqq
\sup \, \bigl\{ L(P,f) : P \text{ partisi dari } [a,b] \bigr\} , \\
& \overline{\int_a^b} f(x)\,dx \coloneqq
\inf \, \bigl\{ U(P,f) : P \text{ partisi dari } [a,b] \bigr\} .
\end{align*}
Kita menyebut $\underline{\int}$\glsadd{not:lowerdarboux}
sebagai \emph{\myindex{integral Darboux bawah}} dan
$\overline{\int}$\glsadd{not:upperdarboux} sebagai
\emph{\myindex{integral Darboux atas}}\index{integral Darboux}.
Agar tidak perlu mempermasalahkan variabel integrasi, 
kita sering kali cukup menulis
\begin{equation*}
\underline{\int_a^b} f \coloneqq
\underline{\int_a^b} f(x)\,dx 
\qquad \text{dan} \qquad
\overline{\int_a^b} f \coloneqq
\overline{\int_a^b} f(x)\,dx  .
\end{equation*}
\end{defn}

Agar pengintegralan masuk akal, integral Darboux bawah dan atas
seharusnya merupakan bilangan yang sama karena kita menginginkan satu bilangan untuk disebut
\emph{integral}.  Namun, kedua integral ini mungkin berbeda untuk
fungsi-fungsi tertentu.

\begin{example} \label{example:dirichletfunc}
Pertimbangkan \myindex{fungsi Dirichlet}
$f \colon [0,1] \to \R$, dengan $f(x) \coloneqq 1$ jika
$x \in \Q$ dan $f(x) \coloneqq 0$ jika $x \notin \Q$.  Maka
\begin{equation*}
\underline{\int_0^1} f = 0 \qquad \text{dan} \qquad
\overline{\int_0^1} f = 1 .
\end{equation*}
Alasannya, untuk setiap partisi $P$ dan setiap $i$ berlaku 
$m_i = \inf \bigl\{ f(x) : x \in [x_{i-1},x_i] \bigr\} = 0$  dan
$M_i = \sup \bigl\{ f(x) : x \in [x_{i-1},x_i] \bigr\} = 1$.  Dengan demikian,
\begin{equation*}
L(P,f) = \sum_{i=1}^n 0 \cdot \Delta x_i = 0 , \quad \text{dan} \quad
U(P,f) = \sum_{i=1}^n 1 \cdot \Delta x_i = \sum_{i=1}^n \Delta x_i = 1  .
\end{equation*}
\end{example}

\begin{remark}
Definisi yang sama untuk $\underline{\int_a^b} f$ dan
$\overline{\int_a^b} f$
digunakan ketika $f$ didefinisikan pada himpunan lebih besar $S$ sedemikian sehingga
$[a,b] \subset S$.  Dalam hal ini, kita menggunakan pembatasan $f$ pada $[a,b]$
dan harus memastikan bahwa pembatasan tersebut terbatas pada $[a,b]$.
\end{remark}

Untuk menghitung integral, kita sering mengambil partisi $P$ lalu memperhalusnya.
Artinya, kita membagi interval-interval dalam partisi menjadi bagian-bagian yang lebih kecil lagi.

\begin{defn}
Misalkan $P = \{ x_0, x_1, \ldots, x_n \}$ dan
$\widetilde{P} = \{ \widetilde{x}_0, \widetilde{x}_1, \ldots,
\widetilde{x}_{\ell} \}$ merupakan
partisi dari $[a,b]$.  Kita mengatakan bahwa $\widetilde{P}$ adalah
\emph{penghalusan}\index{penghalusan partisi} dari $P$
jika, sebagai himpunan, $P \subset \widetilde{P}$.
\end{defn}

Artinya, $\widetilde{P}$ merupakan penghalusan suatu partisi jika memuat semua
bilangan dalam $P$ dan mungkin beberapa bilangan lain di antaranya.  Sebagai contoh,
$\{ 0, 0.5, 1, 2 \}$ merupakan partisi dari $[0,2]$ dan
$\{ 0, 0.2, 0.5, 1, 1.5, 1.75, 2 \}$ merupakan penghalusannya.
Alasan utama memperkenalkan penghalusan adalah proposisi berikut.

\begin{prop} \label{prop:refinement}
Misalkan $f \colon [a,b] \to \R$ merupakan fungsi terbatas dan $P$
merupakan partisi dari $[a,b]$.  Misalkan $\widetilde{P}$ merupakan penghalusan $P$.
Maka
\begin{equation*}
L(P,f) \leq L(\widetilde{P},f) 
\qquad \text{dan} \qquad
U(\widetilde{P},f) \leq U(P,f) .
\end{equation*}
\end{prop}

\begin{proof}
Bagian pelik dari pembuktian ini adalah menuliskan notasinya dengan benar.
Misalkan $\widetilde{P} = \{ \widetilde{x}_0, \widetilde{x}_1, \ldots,
\widetilde{x}_{\ell} \}$ merupakan
penghalusan dari 
$P = \{ x_0, x_1, \ldots, x_n \}$.  Maka
$x_0 = \widetilde{x}_0$ dan 
$x_n = \widetilde{x}_{\ell}$.  Bahkan, terdapat bilangan bulat
$k_0 < k_1 < \cdots < k_n$ sedemikian sehingga $x_i = \widetilde{x}_{k_i}$ untuk
$i=0,1,2,\ldots,n$.

Misalkan $\Delta \widetilde{x}_q \coloneqq \widetilde{x}_q - \widetilde{x}_{q-1}$ untuk
$q=1,2,\ldots,\ell$.
Lihat \figureref{fig:refinement}.
Kita memperoleh 
\begin{equation*}
\Delta x_i
=
x_i - x_{i-1} =
\widetilde{x}_{k_i} - \widetilde{x}_{k_{i-1}} =
\sum_{q=k_{i-1}+1}^{k_i} 
\bigl(\widetilde{x}_{q} - \widetilde{x}_{q-1}\bigr)
=
\sum_{q=k_{i-1}+1}^{k_i} \Delta \widetilde{x}_q .
\end{equation*}
\begin{myfigureht}
\myincludepdft{figrefinement}{%
Diagram sebuah interval antara x subskrip i minus 1 dan x subskrip i dengan
lebar Delta x subskrip i.
Interval tersebut dibagi menjadi tiga interval yang lebih kecil.
Titik ujung kiri x subskrip i minus 1
juga ditandai sebagai x tilde subskrip q minus 3
dan x tilde subskrip k subskrip i minus 1.
Demikian pula, titik ujung kanan x subskrip i
juga ditandai sebagai x tilde subskrip q dan x tilde subskrip k subskrip i.
Kedua titik di bagian dalam ditandai sebagai
x tilde subskrip q minus 2 dan
x tilde subskrip q minus 1.
Ketiga interval tersebut memiliki panjang
Delta x tilde subskrip q minus 2,
Delta x tilde subskrip q minus 1,
dan
Delta x tilde subskrip q.}
\caption{Penghalusan suatu subinterval.  Perhatikan bahwa $\Delta x_i =
\Delta \widetilde{x}_{q-2} +
\Delta \widetilde{x}_{q-1} +
\Delta \widetilde{x}_{q}$,
dan juga
$k_{i-1}+1 = q-2$ dan
$k_{i} = q$.\label{fig:refinement}}
\end{myfigureht}

Seperti sebelumnya, misalkan $m_i$ bersesuaian dengan partisi $P$.
Misalkan $\widetilde{m}_q \coloneqq \inf \bigl\{ f(x) : \widetilde{x}_{q-1} \leq x \leq
\widetilde{x}_q \bigr\}$.
Selanjutnya, $m_i \leq \widetilde{m}_q$ untuk $k_{i-1} < q \leq k_i$.  Oleh karena itu,
\begin{equation*}
m_i \Delta x_i
=
m_i \sum_{q=k_{i-1}+1}^{k_i} \Delta \widetilde{x}_q
=
\sum_{q=k_{i-1}+1}^{k_i} m_i \Delta \widetilde{x}_q
\leq
\sum_{q=k_{i-1}+1}^{k_i} \widetilde{m}_q \Delta \widetilde{x}_q .
\end{equation*}
Dengan demikian,
\begin{equation*}
L(P,f) =
\sum_{i=1}^n m_i \Delta x_i
\leq
\sum_{i=1}^n \,
\sum_{q=k_{i-1}+1}^{k_i} \widetilde{m}_q \Delta \widetilde{x}_q
=
\sum_{q=1}^{\ell}
\widetilde{m}_q \Delta \widetilde{x}_q = L(\widetilde{P},f).
\end{equation*}

Pembuktian $U(\widetilde{P},f) \leq U(P,f)$ diserahkan sebagai latihan.
\end{proof}

Berbekal penghalusan, kita membuktikan hasil berikut.
Pokok penting dari proposisi berikutnya adalah bahwa
integral Darboux bawah kurang dari atau sama dengan integral Darboux
atas.

\begin{prop} \label{intulbound:prop}
Misalkan $f \colon [a,b] \to \R$ merupakan fungsi terbatas.  Misalkan $m, M \in \R$ 
sedemikian sehingga untuk semua $x \in [a,b]$ berlaku $m \leq f(x) \leq M$.  Maka
\begin{equation}
\label{intulbound:eq}
m(b-a) \leq
\underline{\int_a^b} f \leq \overline{\int_a^b} f
\leq M(b-a) .
\end{equation}
\end{prop}

\begin{proof}
Berdasarkan \propref{sumulbound:prop}, untuk setiap partisi $P$,
\begin{equation*}
m(b-a) \leq L(P,f) \leq U(P,f) \leq M(b-a).
\end{equation*}
Ketaksamaan
$m(b-a) \leq L(P,f)$ menyiratkan $m(b-a) \leq \underline{\int_a^b} f$.
Ketaksamaan
$U(P,f) \leq M(b-a)$ menyiratkan $\overline{\int_a^b} f \leq M(b-a)$.

Ketaksamaan tengah dalam
\eqref{intulbound:eq} merupakan pokok utama proposisi ini.
Misalkan $P_1, P_2$ merupakan partisi dari $[a,b]$.  Definisikan 
$\widetilde{P} \coloneqq P_1 \cup P_2$.
Himpunan $\widetilde{P}$ merupakan partisi dari $[a,b]$ yang
merupakan penghalusan dari $P_1$ sekaligus penghalusan dari $P_2$.
Berdasarkan \propref{prop:refinement},
$L(P_1,f) \leq L(\widetilde{P},f)$ dan
$U(\widetilde{P},f) \leq U(P_2,f)$.  Jadi,
\begin{equation*}
L(P_1,f) \leq L(\widetilde{P},f) \leq U(\widetilde{P},f) \leq U(P_2,f) .
\end{equation*}
Dengan kata lain, untuk sebarang dua partisi $P_1$ dan $P_2$, berlaku
$L(P_1,f) \leq U(P_2,f)$.  
Ingat kembali \propref{infsupineq:prop}, lalu ambil supremum dan
infimum terhadap semua partisi:
\begin{multline*}
\underline{\int_a^b} f = 
\sup \, \bigl\{ L(P,f) : P \text{ partisi dari } [a,b] \bigr\}
\\
\leq
\inf \, \bigl\{ U(P,f) : P \text{ partisi dari } [a,b] \bigr\}
=
\overline{\int_a^b} f . \qedhere
\end{multline*}
\end{proof}

\subsection{Integral Riemann}

Akhirnya kita dapat mendefinisikan integral Riemann.  Namun, integral Riemann
hanya didefinisikan pada kelas fungsi tertentu, yang anggotanya disebut
fungsi yang terintegralkan secara Riemann.

\begin{defn}
Misalkan $f \colon [a,b] \to \R$ suatu fungsi terbatas sedemikian sehingga
\begin{equation*}
\underline{\int_a^b} f(x)\,dx = \overline{\int_a^b} f(x)\,dx .
\end{equation*}
Maka $f$ dikatakan \emph{\myindex{terintegralkan secara Riemann}}.
Himpunan fungsi yang terintegralkan secara Riemann pada $[a,b]$ dilambangkan
dengan $\sR\bigl([a,b]\bigr)$.\glsadd{not:integrablefunc}
Ketika $f \in \sR\bigl([a,b]\bigr)$, kita definisikan\glsadd{not:riemannint}
\begin{equation*}
\int_a^b f(x)\,dx \coloneqq 
\underline{\int_a^b} f(x)\,dx = \overline{\int_a^b} f(x)\,dx .
\end{equation*}
Seperti sebelumnya, kita sering menulis
\begin{equation*}
\int_a^b f \coloneqq \int_a^b f(x)\,dx.
\end{equation*}
Bilangan $\int_a^b f$ disebut \emph{\myindex{integral Riemann}}
dari $f$, atau kadang-kadang cukup \emph{integral} dari $f$.
\end{defn}

Berdasarkan definisi, fungsi yang terintegralkan secara Riemann adalah fungsi terbatas.
Dengan menerapkan \propref{intulbound:prop}, kita langsung memperoleh
proposisi berikut.  Lihat juga \figureref{fig:integralminmax}.

\begin{prop} \label{intbound:prop}
Misalkan $f \colon [a,b] \to \R$ suatu fungsi yang terintegralkan secara Riemann.
Misalkan $m, M \in \R$ 
sedemikian sehingga $m \leq f(x) \leq M$ untuk semua $x \in [a,b]$.  Maka
\begin{equation*}
m(b-a) \leq
\int_a^b f
\leq M(b-a) .
\end{equation*}
\end{prop}
\begin{myfigureht}
\myincludegraphics{integralminmax}{%
Diagram sebuah persegi panjang di bidang
yang secara horizontal berada di antara a dan b, dan secara vertikal berada di antara sumbu x
dan M kapital.  Persegi panjang tersebut terbagi menjadi bagian atas
yang berwarna putih dan memuat grafik suatu fungsi, serta bagian bawah
yang diarsir.  Garis horizontal pembatasnya diberi label
m kecil.}
\caption{Luas di bawah kurva dibatasi dari atas oleh
luas seluruh persegi panjang, $M(b-a)$, dan dari bawah oleh
luas bagian yang diarsir, $m(b-a)$.\label{fig:integralminmax}}
\end{myfigureht}

Bentuk yang lebih lemah dari proposisi ini sering kali berguna:
Jika $\babs{f(x)} \leq M$ untuk semua $x \in [a,b]$, maka
\begin{equation*}
\abs{\int_a^b f} \leq M(b-a) .
\end{equation*}

\begin{example}
Kita mengintegralkan fungsi konstan menggunakan
\propref{intulbound:prop}.
Jika $f(x) \coloneqq c$ untuk suatu konstanta $c$, kita ambil $m = M = c$.
Dalam pertidaksamaan \eqref{intulbound:eq},
semua pertidaksamaan tersebut harus berupa kesamaan.
Jadi $f$ terintegralkan pada $[a,b]$ dan $\int_a^b f = c(b-a)$.
\end{example}

\begin{example}
Misalkan $f \colon [0,2] \to \R$ didefinisikan oleh
\begin{equation*}
f(x) \coloneqq
\begin{cases}
1               & \text{jika } x < 1,\\
\nicefrac{1}{2} & \text{jika } x = 1,\\
0               & \text{jika } x > 1.
\end{cases}
\end{equation*}
Kita akan menunjukkan bahwa $f$ terintegralkan secara Riemann dan $\int_0^2 f = 1$.

Bukti: Misalkan $0 < \epsilon < 1$ sebarang.
Misalkan $P \coloneqq \{0, 1-\epsilon, 1+\epsilon, 2\}$ suatu partisi.  Kita menggunakan notasi dari
definisi jumlah Darboux.  Maka
\begin{align*}
m_1 &= \inf  \bigl\{ f(x) : x \in [0,1-\epsilon] \bigr\} = 1 , & 
M_1 &= \sup  \bigl\{ f(x) : x \in [0,1-\epsilon] \bigr\} = 1 , \\
m_2 &= \inf  \bigl\{ f(x) : x \in [1-\epsilon,1+\epsilon] \bigr\} = 0 , & 
M_2 &= \sup  \bigl\{ f(x) : x \in [1-\epsilon,1+\epsilon] \bigr\} = 1 , \\
m_3 &= \inf  \bigl\{ f(x) : x \in [1+\epsilon,2] \bigr\} = 0 , & 
M_3 &= \sup  \bigl\{ f(x) : x \in [1+\epsilon,2] \bigr\} = 0 .
\end{align*}
Selain itu, $\Delta x_1 = 1-\epsilon$, $\Delta x_2 = 2\epsilon$, dan
$\Delta x_3 = 1-\epsilon$.
Lihat \figureref{darbouxfigstep}.
\begin{myfigureht}
\myincludegraphics{darbouxfigstep}{%
Diagram grafik suatu fungsi yang bernilai 1 dari 0 sampai 1,
pada 1 nilainya setengah, dan dari 1 sampai 2 fungsi tersebut bernilai 0.
Sumbu horizontal dibagi menjadi tiga interval,
dari 0 sampai 1 dikurangi epsilon, lalu sampai 1 ditambah epsilon, lalu sampai 2.
Panjangnya masing-masing adalah Delta x subskrip 1 sama dengan 1 dikurangi epsilon,
Delta x subskrip 2 sama dengan 2 epsilon, dan 
Delta x subskrip 3 sama dengan 1 dikurangi epsilon.
Terdapat sebuah persegi panjang berarsir setinggi 1 untuk x antara 0 dan
1 dikurangi epsilon.  Terdapat sebuah persegi panjang putih setinggi 1
untuk x antara 1 dikurangi epsilon dan 1 ditambah epsilon.
Pada sumbu vertikal,
titik bawah ditandai sebagai M subskrip 3 sama dengan m subskrip 2
sama dengan m subskrip 3 sama dengan 0.
Titik atas ditandai sebagai M subskrip 1 sama dengan M subskrip 2
sama dengan m subskrip 1 sama dengan 1.}
\caption{Jumlah Darboux untuk fungsi tangga.  $L(P,f)$ adalah luas persegi panjang
yang diarsir, $U(P,f)$ adalah luas gabungan kedua persegi panjang, dan
$U(P,f)-L(P,f)$ adalah luas persegi panjang yang tidak diarsir.\label{darbouxfigstep}}
\end{myfigureht}

Kita hitung
\begin{align*}
& L(P,f) = \sum_{i=1}^3 m_i \Delta x_i =
1 \cdot (1-\epsilon) + 0 \cdot 2\epsilon + 0 \cdot (1-\epsilon)
= 1-\epsilon , \\
& U(P,f) = \sum_{i=1}^3 M_i \Delta x_i =
1 \cdot (1-\epsilon) + 1 \cdot 2\epsilon + 0 \cdot (1-\epsilon)
= 1+\epsilon .
\end{align*}
Dengan demikian,
\begin{equation*}
\overline{\int_0^2} f - 
\underline{\int_0^2} f
\leq
U(P,f) - L(P,f)
=
(1+\epsilon)
- (1-\epsilon) = 2 \epsilon .
\end{equation*}
Menurut \propref{intulbound:prop}, $\underline{\int_0^2} f \leq \overline{\int_0^2} f$.
Karena $\epsilon$ sebarang,
$\overline{\int_0^2} f = \underline{\int_0^2} f$.  Jadi $f$ terintegralkan secara
Riemann.  Terakhir,
\begin{equation*}
1-\epsilon = L(P,f) \leq \int_0^2 f \leq U(P,f) =
1+\epsilon.
\end{equation*}
Oleh karena itu, $\bigl\lvert \int_0^2 f - 1 \bigr\rvert \leq \epsilon$.  Karena $\epsilon$ sebarang,
kita simpulkan $\int_0^2 f = 1$.
\end{example}

Ada baiknya sebagian teknik dalam contoh tersebut kita sarikan menjadi sebuah
proposisi.  Perhatikan bahwa $U(P,f)-L(P,f)$ tepat sama dengan luas total bagian-bagian
persegi panjang yang tidak diarsir pada \figureref{darbouxfig}.

\begin{prop}
Misalkan $f \colon [a,b] \to \R$ suatu fungsi terbatas.  Maka $f$ terintegralkan secara
Riemann jika untuk setiap $\epsilon > 0$, terdapat suatu partisi $P$ dari
$[a,b]$ sedemikian sehingga
\begin{equation*}
U(P,f) - L(P,f) < \epsilon .
\end{equation*}
\end{prop}

\begin{proof}
Jika untuk setiap $\epsilon > 0$ terdapat $P$ seperti itu, maka
\begin{equation*}
0 \leq
\overline{\int_a^b} f - 
\underline{\int_a^b} f
\leq
U(P,f) - L(P,f) < \epsilon .
\end{equation*}
Oleh karena itu, 
$\overline{\int_a^b} f = \underline{\int_a^b} f$, dan $f$ terintegralkan.
\end{proof}

\begin{example}
Mari kita tunjukkan bahwa $\frac{1}{1+x}$ terintegralkan pada $[0,b]$ untuk setiap $b > 0$.
Kelak kita akan melihat bahwa fungsi kontinu terintegralkan, tetapi di sini kita
membuktikannya secara langsung.

Diberikan $\epsilon > 0$.  Ambil $n \in \N$ dan
tetapkan $x_i \coloneqq \nicefrac{ib}{n}$ untuk membentuk 
partisi $P \coloneqq \{ x_0,x_1,\ldots,x_n \}$ dari $[0,b]$.
Maka $\Delta x_i = \nicefrac{b}{n}$ untuk setiap $i$.  
Karena $f$ menurun, untuk setiap subinterval $[x_{i-1},x_i]$,
\begin{equation*}
%mbxlatex \begin{aligned}
m_i
%mbxlatex &
=
\inf \left\{ \frac{1}{1+x} : x \in [x_{i-1},x_i] \right\} = \frac{1}{1+x_i} ,
%mbxSTARTIGNORE
\quad
%mbxENDIGNORE
%mbxlatex \\
M_i
=
%mbxlatex &
\sup \left\{ \frac{1}{1+x} : x \in [x_{i-1},x_i] \right\} =
\frac{1}{1+x_{i-1}} .
%mbxlatex \end{aligned}
\end{equation*}
Maka
\begin{multline*}
U(P,f)-L(P,f)
=
\sum_{i=1}^n
\Delta x_i
(M_i-m_i)
=
\frac{b}{n}
\sum_{i=1}^n 
\left( \frac{1}{1+\nicefrac{(i-1)b}{n}} - \frac{1}{1+\nicefrac{ib}{n}} \right) 
% Tanda sama dengan redundan pada sumber dihilangkan.
\\
=
\frac{b}{n}
\left( \frac{1}{1+\nicefrac{0b}{n}} - \frac{1}{1+\nicefrac{nb}{n}} \right) 
=
\frac{b^2}{n(b+1)} .
\end{multline*}
Jumlah tersebut bersifat teleskopik---suku-sukunya saling menghapus secara berturut-turut,
sesuatu yang telah kita jumpai sebelumnya.
Dengan memilih $n$ sedemikian sehingga
$\frac{b^2}{n(b+1)} < \epsilon$, syarat dalam proposisi terpenuhi, dan fungsi
tersebut terintegralkan.
\end{example}

\begin{remark}
Salah satu cara memahami integral adalah bahwa integral menjumlahkan (mengintegralkan) banyak
informasi lokal---integral menjumlahkan $f(x)\,dx$ untuk semua $x$.
Leibniz memilih tanda integral berupa huruf S yang dipanjangkan untuk menyatakan penjumlahan.
Berbeda dengan turunan, yang bersifat \myquote{lokal,} 
integral muncul dalam penerapan ketika
kita menginginkan jawaban \myquote{global}: jarak total yang ditempuh, suhu
rata-rata, muatan total, dan sebagainya.
\end{remark}

\subsection{Notasi lebih lanjut}

Jika $f \colon S \to \R$ didefinisikan pada himpunan yang lebih besar $S$ dan
$[a,b] \subset S$,
kita mengatakan bahwa $f$ terintegralkan Riemann pada $[a,b]$ jika pembatasan $f$
ke $[a,b]$ terintegralkan Riemann.
Dalam hal ini,
kita mengatakan $f \in \sR\bigl([a,b]\bigr)$,
dan
kita menuliskan $\int_a^b f$ untuk menyatakan integral Riemann
dari pembatasan $f$ ke $[a,b]$.

Berguna untuk mendefinisikan integral $\int_a^b f$ bahkan jika
$a \nless b$.  Andaikan $b < a$ dan $f \in \sR\bigl([b,a]\bigr)$.
Definisikan
\begin{equation*}
\int_a^b f \coloneqq - \int_b^a f .
\end{equation*}
Untuk sebarang fungsi $f$, definisikan
\begin{equation*}
\int_a^a f \coloneqq 0 .
\end{equation*}

Terkadang, variabel $x$ mungkin sudah memiliki arti lain.  Ketika
kita perlu menuliskan variabel integrasi, kita dapat
menggunakan huruf lain.  Sebagai contoh,
\begin{equation*}
\int_a^b f(s)\,ds \coloneqq \int_a^b f(x)\,dx .
\end{equation*}

\subsection{Latihan}

\begin{exercise}
Definisikan $f \colon [0,1] \to \R$ dengan $f(x) \coloneqq x^3$
dan ambil $P \coloneqq \{ 0, 0.1, 0.4, 1 \}$.  Hitung $L(P,f)$ dan $U(P,f)$.
\end{exercise}

\begin{exercise}
Misalkan $f \colon [0,1] \to \R$ didefinisikan dengan $f(x) \coloneqq x$.
Tunjukkan bahwa $f \in \sR\bigl([0,1]\bigr)$ dan
hitung $\int_0^1 f$ menggunakan definisi integral
(tetapi
Anda boleh menggunakan proposisi-proposisi dalam bagian ini).%\propref{intulbound:prop}).
\end{exercise}

\begin{exercise}
Misalkan $f \colon [a,b] \to \R$ adalah fungsi terbatas.
Andaikan terdapat barisan partisi $\{ P_k \}_{k=1}^\infty$ dari $[a,b]$
sedemikian sehingga
\begin{equation*}
\lim_{k \to \infty} \bigl( U(P_k,f) - L(P_k,f) \bigr) = 0 .
\end{equation*}
Tunjukkan bahwa $f$ terintegralkan Riemann dan bahwa
\begin{equation*}
\int_a^b f = 
\lim_{k \to \infty} U(P_k,f)
=
\lim_{k \to \infty} L(P_k,f) .
\end{equation*}
\end{exercise}

\begin{exercise}
Selesaikan pembuktian \propref{prop:refinement}.
\end{exercise}

\begin{exercise}
Andaikan $f \colon [-1,1] \to \R$ didefinisikan sebagai
\begin{equation*}
f(x) \coloneqq
\begin{cases}
1 & \text{jika } x > 0, \\
0 & \text{jika } x \leq 0.
\end{cases}
\end{equation*}
Buktikan bahwa $f \in \sR\bigl([-1,1]\bigr)$ dan
hitung $\int_{-1}^1 f$ menggunakan definisi integral
(tetapi
Anda boleh menggunakan proposisi-proposisi dalam bagian ini).
%(feel free to use \propref{intulbound:prop}).
\end{exercise}

\begin{exercise}
Misalkan $c \in (a,b)$ dan $d \in \R$.
Definisikan $f \colon [a,b] \to \R$ sebagai
\begin{equation*}
f(x) \coloneqq
\begin{cases}
d & \text{jika } x = c, \\
0 & \text{jika } x \neq c.
\end{cases}
\end{equation*}
Buktikan bahwa $f \in \sR\bigl([a,b]\bigr)$ dan
hitung
$\int_a^b f$ menggunakan definisi integral
%(feel free to use \propref{intulbound:prop}).
(tetapi
Anda boleh menggunakan proposisi-proposisi dalam bagian ini).
\end{exercise}

\begin{exercise} \label{exercise:taggedpartition}
Andaikan $f \colon [a,b] \to \R$ terintegralkan Riemann.  Diberikan $\epsilon
> 0$.  Tunjukkan bahwa terdapat partisi $P = \{ x_0, x_1,
\ldots, x_n \}$
sedemikian sehingga untuk setiap
himpunan bilangan $\{ c_1, c_2, \ldots, c_n \}$ dengan
$c_k \in [x_{k-1},x_k]$ untuk setiap $k$, berlaku
\begin{equation*}
\abs{\int_a^b f - \sum_{k=1}^n f(c_k) \Delta x_k} < \epsilon .
\end{equation*}
\end{exercise}

\begin{exercise}
Misalkan $f \colon [a,b] \to \R$ adalah fungsi yang terintegralkan Riemann.
Misalkan $\alpha > 0$ dan $\beta \in \R$.
Definisikan $g(x) \coloneqq f(\alpha x + \beta)$ pada interval
$I = [\frac{a-\beta}{\alpha}, \frac{b-\beta}{\alpha}]$.  Tunjukkan
bahwa $g$ terintegralkan Riemann pada $I$.
\end{exercise}

\begin{exercise}
Andaikan $f \colon [0,1] \to \R$ dan $g \colon [0,1] \to \R$
sedemikian sehingga untuk setiap $x \in (0,1]$,
berlaku $f(x) = g(x)$.  Andaikan $f$ terintegralkan Riemann. 
Buktikan bahwa $g$ terintegralkan Riemann dan $\int_{0}^1 f = \int_{0}^1 g$.
\end{exercise}

\begin{exercise}
Misalkan $f \colon [0,1] \to \R$ adalah fungsi terbatas.
Misalkan $P_n = \{ x_0,x_1,\ldots,x_n \}$ adalah partisi seragam dari $[0,1]$,
yaitu $x_i = \nicefrac{i}{n}$.  Apakah $\bigl\{ L(P_n,f) \bigr\}_{n=1}^\infty$
selalu monoton?  Ya/Tidak: Buktikan atau temukan contoh penyangkal.
\end{exercise}

\begin{exercise}[Menantang]
Untuk fungsi terbatas $f \colon [0,1] \to \R$, misalkan
$R_n \coloneqq (\nicefrac{1}{n})\sum_{i=1}^n f(\nicefrac{i}{n})$ (aturan
titik ujung kanan seragam).
\begin{enumerate}[a)]
\item
Jika $f$ terintegralkan Riemann, tunjukkan bahwa $\int_0^1 f = \lim\limits_{n\to\infty} R_n$.
\item
Temukan $f$ yang tidak terintegralkan Riemann, tetapi $\lim\limits_{n\to\infty} R_n$ ada.
\end{enumerate}
\end{exercise}

\begin{exercise}[Menantang] \label{exercise:riemannintdarboux}
Generalisasikan latihan sebelumnya.
Tunjukkan bahwa $f \in \sR\bigl([a,b]\bigr)$ jika dan hanya jika terdapat $I \in \R$,
sedemikian sehingga untuk setiap $\epsilon > 0$ terdapat
$\delta > 0$ sedemikian sehingga jika $P$ adalah partisi dengan $\Delta x_i < \delta$
untuk setiap $i$, maka
$\babs{L(P,f) - I} < \epsilon$ dan
$\babs{U(P,f) - I} < \epsilon$.  Jika $f \in \sR\bigl([a,b]\bigr)$, maka $I = \int_a^b f$.
\end{exercise}

\begin{exercise}
Dengan menggunakan \exerciseref{exercise:riemannintdarboux} dan gagasan
pembuktian dalam \exerciseref{exercise:taggedpartition}, tunjukkan bahwa 
integral Darboux ekuivalen dengan definisi standar
integral Riemann, yang kemungkinan besar telah Anda jumpai dalam kalkulus.  Artinya,
tunjukkan bahwa
$f \in \sR\bigl([a,b]\bigr)$ jika dan hanya jika terdapat $I \in \R$,
sedemikian sehingga untuk setiap $\epsilon > 0$ terdapat
$\delta > 0$ sedemikian sehingga jika $P = \{ x_0,x_1,\ldots,x_n \}$
adalah partisi dengan $\Delta x_i < \delta$
untuk setiap $i$, maka
$\abs{\sum_{i=1}^n f(c_i) \Delta x_i - I} < \epsilon$ untuk setiap himpunan
$\{ c_1,c_2,\ldots,c_n \}$ dengan $c_i \in [x_{i-1},x_i]$.
Jika $f \in \sR\bigl([a,b]\bigr)$, maka $I = \int_a^b f$.
\end{exercise}


\begin{exercise}[Menantang]
Konstruksikan fungsi $f$ dan $g$, 
dengan
$f \colon [0,1] \to \R$ terintegralkan Riemann,
$g \colon [0,1] \to [0,1]$ injektif dan surjektif,
dan sedemikian sehingga komposisi $f \circ g$ tidak terintegralkan Riemann.
\end{exercise}

\begin{exercise}
Andaikan bahwa $f \colon [a,b] \to \R$ adalah fungsi terbatas, dan $P$ adalah
partisi dari $[a,b]$ sedemikian sehingga $L(P,f)=U(P,f)$.  Buktikan bahwa $f$ adalah
fungsi konstan.
\end{exercise}


%%%%%%%%%%%%%%%%%%%%%%%%%%%%%%%%%%%%%%%%%%%%%%%%%%%%%%%%%%%%%%%%%%%%%%%%%%%%%%

\sectionnewpage
\section{Sifat-sifat integral}
\label{sec:rintprop}

%mbxINTROSUBSECTION

\sectionnotes{2 kuliah; keterintegralan fungsi dengan
diskontinuitas dapat dilewati dengan aman}

\subsection{Aditivitas}

Menjumlahkan sekumpulan objek dalam dua bagian lalu menjumlahkan kedua bagian tersebut
seharusnya sama dengan menjumlahkan semuanya sekaligus.
Sifat yang bersesuaian untuk integral disebut
\myindex{sifat aditif integral}.  Pertama-tama, kita membuktikan sifat aditif
untuk integral Darboux bawah dan atas.

\begin{lemma} \label{lemma:darbouxadd}
Andaikan $a < b < c$ dan $f \colon [a,c] \to \R$ fungsi terbatas.
Maka
\begin{equation*}
\underline{\int_a^c} f
=
\underline{\int_a^b} f
+
\underline{\int_b^c} f
\quad \text{dan} \quad
\overline{\int_a^c} f
=
\overline{\int_a^b} f
+
\overline{\int_b^c} f .
\end{equation*}
\end{lemma}

\begin{proof}
Jika kita memiliki partisi $P_1 = \{ x_0,x_1,\ldots,x_k \}$
dari $[a,b]$ dan $P_2 = \{ x_k, x_{k+1}, \ldots, x_n \}$ dari $[b,c]$,
maka himpunan $P \coloneqq P_1 \cup P_2 = \{ x_0, x_1, \ldots, x_n \}$ merupakan
partisi dari $[a,c]$.  Kita peroleh
\begin{equation*}
L(P,f) =
\sum_{i=1}^n m_i \Delta x_i
=
\sum_{i=1}^k m_i \Delta x_i
+
\sum_{i=k+1}^n m_i \Delta x_i
=
L(P_1,f) + L(P_2,f) .
\end{equation*}
Ketika mengambil supremum ruas kanan atas semua $P_1$ dan $P_2$,
kita mengambil supremum ruas kiri
atas semua partisi $P$ dari $[a,c]$ yang memuat $b$.  Jika $Q$ merupakan partisi
dari $[a,c]$ dan $P = Q \cup \{ b \}$, maka $P$ merupakan penghalusan dari $Q$
sehingga $L(Q,f) \leq L(P,f)$.  Oleh karena itu, mengambil supremum hanya atas $P$
yang memuat $b$ sudah cukup untuk memperoleh supremum $L(P,f)$
atas semua partisi $P$; lihat \exerciseref{exercise:dominatingb}.
Terakhir, ingat kembali \exerciseref{exercise:supofsum}
untuk menghitung
\begin{equation*}
\begin{split}
\underline{\int_a^c} f
& =
\sup \, \bigl\{ L(P,f) : P \text{ partisi dari } [a,c] \bigr\}
\\
& =
\sup \, \bigl\{ L(P,f) : P \text{ partisi dari } [a,c], b \in P \bigr\}
\\
& =
\sup \, \bigl\{ L(P_1,f) + L(P_2,f) :
P_1 \text{ partisi dari } [a,b], P_2 \text{ partisi dari } [b,c] \bigr\}
\\
& =
\sup \, \bigl\{ L(P_1,f) : P_1 \text{ partisi dari } [a,b] \bigr\}
+
\sup \, \bigl\{ L(P_2,f) : P_2 \text{ partisi dari } [b,c] \bigr\}
\\
&=
\underline{\int_a^b} f + \underline{\int_b^c} f .
\end{split}
\end{equation*}

Demikian pula, untuk $P$, $P_1$, dan $P_2$ seperti di atas, kita peroleh
\begin{equation*}
U(P,f) =
\sum_{i=1}^n M_i \Delta x_i
=
\sum_{i=1}^k M_i \Delta x_i
+
\sum_{i=k+1}^n M_i \Delta x_i
=
U(P_1,f) + U(P_2,f) .
\end{equation*}
Kita hendak mengambil infimum ruas kanan
atas semua $P_1$ dan $P_2$, sehingga kita mengambil infimum
atas semua partisi $P$ dari $[a,c]$ yang memuat $b$.  Jika $Q$ merupakan partisi
dari $[a,c]$ dan $P = Q \cup \{ b \}$, maka $P$ merupakan penghalusan dari $Q$
sehingga $U(Q,f) \geq U(P,f)$.  Oleh karena itu, mengambil infimum hanya atas $P$
yang memuat $b$ sudah cukup untuk memperoleh infimum $U(P,f)$ atas
semua $P$.
Kita peroleh
\begin{equation*}
\overline{\int_a^c} f
=
\overline{\int_a^b} f + \overline{\int_b^c} f .  \qedhere
\end{equation*}
\end{proof}

\begin{prop}
Misalkan $a < b < c$.  Fungsi $f \colon [a,c] \to \R$ terintegralkan secara Riemann
jika dan hanya jika $f$ terintegralkan secara Riemann pada $[a,b]$ dan $[b,c]$.  Jika
$f$ terintegralkan secara Riemann, maka
\begin{equation*}
\int_a^c f
=
\int_a^b f
+
\int_b^c f .
\end{equation*}
\end{prop}

\begin{proof}
Andaikan $f \in \sR\bigl([a,c]\bigr)$.  Maka fungsi tersebut terbatas dan
$\overline{\int_a^c} f = 
\underline{\int_a^c} f = 
\int_a^c f$.  Lema tersebut memberikan
\begin{equation*}
\int_a^c f
=
\underline{\int_a^c} f
 =
\underline{\int_a^b} f + \underline{\int_b^c} f
 \leq
\overline{\int_a^b} f + \overline{\int_b^c} f
 =
\overline{\int_a^c} f
 =
\int_a^c f .
\end{equation*}
Jadi, pertidaksamaan tersebut berlaku sebagai kesamaan:
\begin{equation*}
\underline{\int_a^b} f + \underline{\int_b^c} f
=
\overline{\int_a^b} f + \overline{\int_b^c} f .
\end{equation*}
Karena kita juga mengetahui
$\underline{\int_a^b} f \leq \overline{\int_a^b} f$
dan
$\underline{\int_b^c} f \leq \overline{\int_b^c} f$, kita
menyimpulkan
\begin{equation*}
\underline{\int_a^b} f
=
\overline{\int_a^b} f
\qquad \text{dan} \qquad
\underline{\int_b^c} f
=
\overline{\int_b^c} f .
\end{equation*}
Jadi, $f$ terintegralkan secara Riemann pada $[a,b]$ dan $[b,c]$, serta rumus yang diinginkan
berlaku.

Sekarang andaikan $f$ terintegralkan secara Riemann pada $[a,b]$ dan pada $[b,c]$.
Sekali lagi, fungsi tersebut terbatas, dan lema tersebut memberikan
\begin{equation*}
\underline{\int_a^c} f
=
\underline{\int_a^b} f + \underline{\int_b^c} f
=
\int_a^b f + \int_b^c f
=
\overline{\int_a^b} f + \overline{\int_b^c} f
=
\overline{\int_a^c} f .
\end{equation*}
Oleh karena itu, $f$ terintegralkan secara Riemann pada $[a,c]$, dan integralnya dihitung
seperti yang dinyatakan.
\end{proof}

Akibat sederhana dari aditivitas adalah korolari berikut.  Kita
menyerahkan perinciannya kepada pembaca sebagai latihan.

\begin{cor} \label{intsubcor}
Jika $f \in \sR\bigl([a,b]\bigr)$ dan
$[c,d] \subset [a,b]$, maka
pembatasan $f|_{[c,d]}$ berada dalam $\sR\bigl([c,d]\bigr)$.
\end{cor}

\subsection{Linearitas dan kemonotonan}

Suatu jumlah merupakan fungsi linear dari suku-sukunya.  Demikian pula integral.

\begin{prop}[Linearitas]
\index{linearitas integral}\label{prop:integrallinear}
Misalkan $f$ dan $g$ berada dalam $\sR\bigl([a,b]\bigr)$ dan $\alpha \in \R$.
\begin{enumerate}[(i)]
\item $\alpha f$ berada dalam $\sR\bigl([a,b]\bigr)$ dan
\begin{equation*}
\int_a^b \alpha f(x) \,dx = \alpha \int_a^b f(x) \,dx .
\end{equation*}
\item $f+g$ berada dalam $\sR\bigl([a,b]\bigr)$ dan
\begin{equation*}
\int_a^b \bigl( f(x)+g(x) \bigr) \,dx = 
\int_a^b f(x) \,dx 
+
\int_a^b g(x) \,dx .
\end{equation*}
\end{enumerate}
\end{prop}

\begin{proof}
\pagebreak[2]
Mari kita buktikan butir pertama untuk $\alpha \geq 0$.
Misalkan $P$ partisi dari $[a,b]$, dan seperti biasa,
$m_i \coloneqq \inf \bigl\{ f(x) : x \in [x_{i-1},x_i] \bigr\}$.
Karena $\alpha \geq 0$, faktor $\alpha$ dapat dikeluarkan
dari infimum,
\begin{equation*}
\inf \bigl\{ \alpha f(x) : x \in [x_{i-1},x_i] \bigr\}
=
\alpha \inf \bigl\{ f(x) : x \in [x_{i-1},x_i] \bigr\} = \alpha m_i .
\end{equation*}
Oleh karena itu,
\begin{equation*}
L(P,\alpha f) =
\sum_{i=1}^n \alpha m_i \Delta x_i = \alpha \sum_{i=1}^n m_i \Delta x_i = \alpha
L(P,f).
\end{equation*}
Demikian pula,
\begin{equation*}
U(P,\alpha f) = \alpha U(P,f) .
\end{equation*}
Sekali lagi, karena $\alpha \geq 0$, kita
dapat mengeluarkan faktor $\alpha$ dari supremum.  Jadi,
\begin{equation*}
\begin{split}
\underline{\int_a^b} \alpha f(x)\,dx & =
\sup \, \bigl\{ L(P,\alpha f) : P \text{ partisi dari } [a,b] \bigr\}
\\
& =
\sup \, \bigl\{ \alpha L(P,f) : P \text{ partisi dari } [a,b] \bigr\}
\\
& =
\alpha \,
\sup \, \bigl\{ L(P,f) : P \text{ partisi dari } [a,b] \bigr\}
\\
& =
\alpha
\underline{\int_a^b} f(x)\,dx .
\end{split}
\end{equation*}
Demikian pula, kita tunjukkan
\begin{equation*}
\overline{\int_a^b} \alpha f(x)\,dx
=
\alpha
\overline{\int_a^b} f(x)\,dx .
\end{equation*}
Kesimpulannya kini berlaku untuk $\alpha \geq 0$.

Untuk menyelesaikan bukti butir pertama (untuk $\alpha < 0$), kita perlu menunjukkan
bahwa $-f$ terintegralkan secara Riemann dan
$\int_a^b - f(x)\,dx =
-
\int_a^b f(x)\,dx$.  Bukti fakta ini diserahkan sebagai
\exerciseref{exercise:proofoflinpropparti}.

Bukti butir kedua diserahkan sebagai
\exerciseref{exercise:proofoflinproppartii}.
Bukti tersebut tidak sulit, tetapi juga tidak sesederhana
yang mungkin tampak pada pandangan pertama.
\end{proof}

Butir kedua dalam proposisi tersebut tidak berlaku sebagai
persamaan untuk integral Darboux, tetapi kita memperoleh pertidaksamaan.
Bukti proposisi berikut merupakan \exerciseref{exercise:upperlowerlinineq}.
Untuk jumlah atas dan bawah pada partisi tetap, hasil ini mengikuti dari \exerciseref{exercise:sumofsup},
yaitu supremum suatu jumlah kurang dari atau sama dengan jumlah
supremumnya, dan demikian pula untuk infimum.

\begin{prop} \label{prop:upperlowerlinineq}
Misalkan $f \colon [a,b] \to \R$ dan $g \colon [a,b] \to \R$ fungsi
terbatas.  Maka
\begin{equation*}
%\overline{\int_a^b} \bigl(f(x)+g(x)\bigr)\,dx \leq
%\overline{\int_a^b}f(x)\,dx+\overline{\int_a^b}g(x)\,dx
\overline{\int_a^b} (f+g) \leq \overline{\int_a^b}f+\overline{\int_a^b}g
,
\qquad
\text{dan}
\qquad
\underline{\int_a^b} (f+g) \geq \underline{\int_a^b}f+\underline{\int_a^b}g
%\underline{\int_a^b} \bigl(f(x)+g(x)\bigr)\,dx \geq
%\underline{\int_a^b}f(x)\,dx+\underline{\int_a^b}g(x)\,dx
.
\end{equation*}
\end{prop}

Menjumlahkan bilangan-bilangan yang lebih kecil seharusnya menghasilkan nilai yang lebih kecil.
Hal ini juga berlaku untuk integral.

\begin{prop}[Kemonotonan]
\index{kemonotonan integral}
Misalkan $f \colon [a,b] \to \R$ dan $g \colon [a,b] \to \R$ fungsi
terbatas, dan $f(x) \leq g(x)$
untuk semua $x \in [a,b]$.  Maka
\begin{equation*}
\underline{\int_a^b} f 
\leq
\underline{\int_a^b} g 
\qquad \text{dan} \qquad
\overline{\int_a^b} f 
\leq
\overline{\int_a^b} g .
\end{equation*}
Selain itu, jika $f$ dan $g$ berada dalam $\sR\bigl([a,b]\bigr)$, maka
\begin{equation*}
\int_a^b f 
\leq
\int_a^b g .
\end{equation*}
\end{prop}

\begin{proof}
Misalkan $P = \{ x_0, x_1, \ldots, x_n \}$ partisi dari $[a,b]$.  Kemudian,
definisikan
\begin{equation*}
m_i \coloneqq \inf \, \bigl\{ f(x) : x \in [x_{i-1},x_i] \bigr\}
\qquad \text{dan} \qquad
\widetilde{m}_i \coloneqq \inf \, \bigl\{ g(x) : x \in [x_{i-1},x_i] \bigr\} .
\end{equation*}
Karena $f(x) \leq g(x)$, kita memiliki $m_i \leq \widetilde{m}_i$.
Oleh karena itu,
\begin{equation*}
L(P,f)
=
\sum_{i=1}^n m_i \Delta x_i
\leq
\sum_{i=1}^n \widetilde{m}_i \Delta x_i
=
L(P,g) .
\end{equation*}
Kita mengambil supremum atas semua $P$ (lihat \propref{prop:funcsupinf}) untuk memperoleh
\begin{equation*}
\underline{\int_a^b} f 
\leq
\underline{\int_a^b} g .
\end{equation*}
Demikian pula, kita memperoleh kesimpulan yang sama untuk integral atas.
Terakhir, jika $f$ dan $g$ terintegralkan secara Riemann, integral bawah dan atas
untuk setiap fungsi berimpit, sehingga kesimpulan tersebut berlaku.
\end{proof}

\subsection{Fungsi kontinu}

Mari kita tunjukkan bahwa fungsi kontinu terintegralkan Riemann.
Kita bahkan dapat mengizinkan beberapa diskontinuitas.  Kita mulai dengan suatu
fungsi yang kontinu pada seluruh interval tertutup $[a,b]$.

\begin{lemma} \label{lemma:contint}
Jika $f \colon [a,b] \to \R$ merupakan fungsi kontinu,
maka $f \in \sR\bigl([a,b]\bigr)$.
\end{lemma}

\begin{proof}
Karena $f$ kontinu pada interval tertutup dan terbatas, fungsi tersebut
terbatas dan
kontinu seragam.
Diberikan $\epsilon > 0$, pilihlah $\delta > 0$ sedemikian sehingga
$\sabs{x-y} < \delta$ mengakibatkan $\babs{f(x)-f(y)} < \frac{\epsilon}{b-a}$.

Misalkan $P = \{ x_0, x_1, \ldots, x_n \}$
merupakan partisi dari $[a,b]$ sedemikian sehingga $\Delta x_i < \delta$ untuk setiap $i = 1,2,
\ldots, n$.  Sebagai contoh,
ambillah $n$ sedemikian sehingga $\frac{b-a}{n} < \delta$, dan
misalkan $x_i \coloneqq \frac{i}{n}(b-a) + a$.
Maka untuk setiap $x, y \in [x_{i-1},x_i]$, kita peroleh 
$\sabs{x-y} \leq \Delta x_i < \delta$, sehingga
\begin{equation*}
f(x)-f(y) \leq \babs{f(x)-f(y)} < \frac{\epsilon}{b-a} .
\end{equation*}
Karena $f$ kontinu pada $[x_{i-1},x_i]$, fungsi tersebut mencapai maksimum dan minimum
pada interval ini.
Misalkan $x$ merupakan titik tempat $f$ mencapai maksimum dan $y$ merupakan titik
tempat $f$ mencapai minimum.  Maka $f(x) = M_i$
dan $f(y) = m_i$ dalam notasi dari definisi integral.
Oleh karena itu,
\begin{equation*}
M_i-m_i = f(x)-f(y) < 
\frac{\epsilon}{b-a} .
\end{equation*}
Dengan demikian,
\begin{equation*}
\begin{split}
\overline{\int_a^b} f - 
\underline{\int_a^b} f 
& \leq
U(P,f) - L(P,f)
\\
& =
\left(
\sum_{i=1}^n
M_i \Delta x_i
\right)
-
\left(
\sum_{i=1}^n
m_i \Delta x_i
\right)
\\
& =
\sum_{i=1}^n
(M_i-m_i) \Delta x_i
\\
& <
\frac{\epsilon}{b-a}
\sum_{i=1}^n
\Delta x_i
\\
& =
\frac{\epsilon}{b-a} (b-a)
= \epsilon .
\end{split}
\end{equation*}
Karena $\epsilon > 0$ dipilih secara sebarang,
\begin{equation*}
\overline{\int_a^b} f = \underline{\int_a^b} f ,
\end{equation*}
dan $f$ terintegralkan Riemann pada $[a,b]$.
\end{proof}

Lema kedua mengatakan bahwa fungsi tersebut hanya perlu \myquote{terintegralkan Riemann
di bagian dalam interval,} asalkan fungsi itu terbatas.  Lema tersebut juga menunjukkan cara
menghitung integralnya.

\begin{lemma} \label{lemma:boundedimpriemann}
Misalkan $f \colon [a,b] \to \R$ merupakan fungsi terbatas, serta
$\{ a_n \}_{n=1}^\infty$ dan $\{b_n \}_{n=1}^\infty$ merupakan barisan sedemikian sehingga
$a < a_n < b_n < b$ untuk setiap $n$, dengan
$\lim_{n\to\infty} a_n = a$ dan $\lim_{n\to\infty} b_n = b$.
Andaikan $f \in \sR\bigl([a_n,b_n]\bigr)$ untuk setiap $n$.
Maka $f \in \sR\bigl([a,b]\bigr)$ dan
\begin{equation*}
\int_a^b f = 
\lim_{n \to \infty} \int_{a_n}^{b_n} f .
\end{equation*}
\end{lemma}

\begin{proof}
Misalkan $M > 0$ merupakan bilangan real sedemikian sehingga $\babs{f(x)} \leq M$.
Karena $(b-a) \geq (b_n-a_n)$,
\begin{equation*}
-M(b-a) \leq
-M(b_n-a_n) \leq
\int_{a_n}^{b_n} f
\leq
M(b_n-a_n) \leq
M(b-a) .
\end{equation*}
Oleh karena itu, barisan bilangan
$\bigl\{ \int_{a_n}^{b_n} f \bigr\}_{n=1}^\infty$ terbatas dan berdasarkan
\hyperref[thm:bwseq]{Bolzano--Weierstrass}
memiliki subbarisan konvergen yang diindeks oleh $n_k$.  Mari kita sebut
$L$ sebagai limit dari subbarisan
$\bigl\{ \int_{a_{n_k}}^{b_{n_k}} f \bigr\}_{k=1}^\infty$.

\lemmaref{lemma:darbouxadd} menyatakan bahwa
integral bawah dan integral atas bersifat aditif,
sedangkan hipotesis menyatakan bahwa
$f$ terintegralkan pada $[a_{n_k},b_{n_k}]$.
Oleh karena itu,
\begin{equation*}
\underline{\int_a^b} f
=
\underline{\int_a^{a_{n_k}}} f
+
\int_{a_{n_k}}^{b_{n_k}} f
+
\underline{\int_{b_{n_k}}^b} f
\geq
-M(a_{n_k}-a)
+
\int_{a_{n_k}}^{b_{n_k}} f
-
M(b-b_{n_k}) .
\end{equation*}
Kita ambil limit ketika $k$ menuju $\infty$ pada ruas kanan,
\begin{equation*}
\underline{\int_a^b} f
\geq
-M\cdot 0
+
L
-
M\cdot 0
= L .
\end{equation*}

Berikutnya kita gunakan sifat aditif dari integral atas,
\begin{equation*}
\overline{\int_a^b} f
=
\overline{\int_a^{a_{n_k}}} f
+
\int_{a_{n_k}}^{b_{n_k}} f
+
\overline{\int_{b_{n_k}}^b} f
\leq
M(a_{n_k}-a)
+
\int_{a_{n_k}}^{b_{n_k}} f
+
M(b-b_{n_k}) .
\end{equation*}
Kita ambil subbarisan yang sama 
$\{ \int_{a_{n_k}}^{b_{n_k}} f \}_{k=1}^\infty$ lalu kita ambil limit 
untuk memperoleh
\begin{equation*}
\overline{\int_a^b} f
\leq
M\cdot 0
+
L
+
M\cdot 0
= L .
\end{equation*}
Jadi $\overline{\int_a^b} f = \underline{\int_a^b} f = L$,
sehingga $f$ terintegralkan Riemann dan $\int_a^b f = L$.
Secara khusus, apa pun
subbarisan konvergen yang kita pilih,
$L$ merupakan bilangan yang sama.

Untuk membuktikan pernyataan terakhir lema tersebut, kita gunakan 
\propref{seqconvsubseqconv:prop}.  Kita telah menunjukkan bahwa setiap
subbarisan konvergen
$\bigl\{ \int_{a_{n_k}}^{b_{n_k}} f \bigr\}_{k=1}^\infty$ konvergen ke $L = \int_a^b f$.
Oleh karena itu, barisan
$\bigl\{ \int_{a_n}^{b_n} f \bigr\}_{n=1}^\infty$ konvergen dengan limit $\int_a^b f$.
\end{proof}

Kita katakan bahwa fungsi $f \colon [a,b] \to \R$ memiliki \emph{\myindex{berhingga banyak
diskontinuitas}} jika terdapat himpunan berhingga $S = \{ x_1, x_2, \ldots, x_n \}
\subset [a,b]$, dan $f$ kontinu
di setiap titik dari $[a,b] \setminus S$.

\begin{thm}
Misalkan $f \colon [a,b] \to \R$ merupakan fungsi terbatas dengan berhingga
banyak diskontinuitas.  Maka $f \in \sR\bigl([a,b]\bigr)$.
\end{thm}

\begin{proof}
Kita membagi interval tersebut menjadi berhingga banyak interval $[a_i,b_i]$
sedemikian sehingga $f$ kontinu
pada bagian dalam $(a_i,b_i)$.  Jika $f$ kontinu pada $(a_i,b_i)$,
maka fungsi tersebut kontinu dan karenanya terintegralkan pada $[c_i,d_i]$ bilamana $a_i < c_i < d_i < b_i$.  Berdasarkan
\lemmaref{lemma:boundedimpriemann},
pembatasan
$f$ pada $[a_i,b_i]$ terintegralkan.
Berdasarkan sifat aditif integral (serta \hyperref[induction:thm]{induksi}),
$f$ terintegralkan pada gabungan interval-interval tersebut.
\end{proof}

\subsection{Lebih lanjut tentang fungsi terintegralkan}

Terkadang lebih mudah (atau perlu)
mengubah nilai tertentu dari suatu fungsi lalu
mengintegralkannya.  Hasil berikut menyatakan
bahwa jika kita mengubah nilai di berhingga
banyak titik, integralnya tidak berubah.

\begin{prop}
Misalkan $f \colon [a,b] \to \R$ terintegralkan secara Riemann.  Misalkan $g \colon [a,b] \to
\R$ sedemikian sehingga $f(x) = g(x)$ untuk semua $x \in [a,b] \setminus S$,
dengan $S$ suatu himpunan hingga.
Maka $g$ terintegralkan secara Riemann dan
\begin{equation*}
\int_a^b g = \int_a^b f.
\end{equation*}
\end{prop}

\begin{proof}[Sketsa bukti]
Dengan menggunakan aditivitas integral, bagi interval $[a,b]$ menjadi
interval-interval yang lebih kecil sedemikian sehingga $f(x) = g(x)$ berlaku untuk semua $x$ kecuali di
titik-titik ujung (perinciannya diserahkan kepada pembaca).

Oleh karena itu, tanpa mengurangi keumuman, andaikan $f(x) = g(x)$ untuk
semua $x \in (a,b)$.  Bukti mengikuti dari \lemmaref{lemma:boundedimpriemann},
dan diserahkan sebagai \exerciseref{exercise:changeendpointsintegral}.
\end{proof}

Terakhir, fungsi monoton (naik atau turun) selalu
terintegralkan secara Riemann.  Buktinya diserahkan kepada pembaca sebagai bagian dari
\exerciseref{exercise:boundedvariationintegrable}.

\begin{prop} \label{prop:monotoneintegrable}
Misalkan $f \colon [a,b] \to \R$ fungsi monoton.  Maka $f \in
\sR\bigl([a,b]\bigr)$.
\end{prop}

\subsection{Latihan}

\begin{exercise} \label{exercise:proofoflinpropparti}
Selesaikan bukti bagian pertama dari \propref{prop:integrallinear}.
Misalkan $f$ merupakan anggota $\sR\bigl([a,b]\bigr)$.  Buktikan bahwa
$-f$ merupakan anggota $\sR\bigl([a,b]\bigr)$ dan 
\begin{equation*}
\int_a^b - f(x) \,dx = - \int_a^b f(x) \,dx .
\end{equation*}
\end{exercise}

\begin{exercise} \label{exercise:proofoflinproppartii}
Buktikan bagian kedua dari \propref{prop:integrallinear}.
Misalkan $f$ dan $g$ anggota $\sR\bigl([a,b]\bigr)$.
Buktikan, tanpa menggunakan \propref{prop:upperlowerlinineq}, bahwa $f+g$ anggota
$\sR\bigl([a,b]\bigr)$ dan
\begin{equation*}
\int_a^b \bigl( f(x)+g(x) \bigr) \,dx = 
\int_a^b f(x) \,dx 
+
\int_a^b g(x) \,dx .
\end{equation*}
Petunjuk: Salah satu caranya adalah menggunakan \propref{prop:refinement} untuk menemukan sebuah partisi tunggal $P$
sedemikian sehingga $U(P,f)-L(P,f) < \nicefrac{\epsilon}{2}$ dan
$U(P,g)-L(P,g) < \nicefrac{\epsilon}{2}$.
\end{exercise}

\begin{exercise} \label{exercise:changeendpointsintegral}
Misalkan $f \colon [a,b] \to \R$ terintegralkan secara Riemann, dan misalkan $g \colon [a,b] \to \R$
sedemikian sehingga $f(x) = g(x)$ untuk semua $x \in (a,b)$.
Buktikan bahwa $g$ terintegralkan secara Riemann dan bahwa
\begin{equation*}
\int_a^b g = \int_a^b f.
\end{equation*}
\end{exercise}

\begin{exercise}
Buktikan \emph{\myindex{teorema nilai rata-rata untuk integral}}:
Jika $f \colon [a,b] \to \R$ kontinu, maka terdapat
$c \in [a,b]$ sedemikian sehingga $\int_a^b f = f(c)(b-a)$.
\end{exercise}

\begin{exercise}
Misalkan $f \colon [a,b] \to \R$ fungsi kontinu sedemikian sehingga $f(x) \geq 0$
untuk semua $x \in [a,b]$ dan $\int_a^b f = 0$.  Buktikan bahwa $f(x) = 0$
untuk semua $x$.
\end{exercise}

\begin{exercise}
Misalkan $f \colon [a,b] \to \R$ fungsi kontinu
dan $\int_a^b f = 0$.  Buktikan bahwa
terdapat $c \in [a,b]$ sedemikian sehingga $f(c) = 0$. (Bandingkan dengan
latihan sebelumnya.)
\end{exercise}

\begin{exercise}
Misalkan $f \colon [a,b] \to \R$ dan $g \colon [a,b] \to \R$
fungsi-fungsi kontinu sedemikian sehingga $\int_a^b f = \int_a^b g$.
Tunjukkan bahwa terdapat $c \in [a,b]$ sedemikian sehingga $f(c) = g(c)$.
\end{exercise}

\begin{exercise}
Misalkan $f \in \sR\bigl([a,b]\bigr)$.  Misalkan $\alpha, \beta, \gamma$ bilangan-bilangan sebarang dalam
$[a,b]$ (tidak harus terurut dengan cara apa pun).  Buktikan 
\begin{equation*}
\int_\alpha^\gamma f =
\int_\alpha^\beta f +
\int_\beta^\gamma f .
\end{equation*}
Ingat kembali arti $\int_a^b f$ jika $b \leq a$.
\end{exercise}

\begin{exercise}
Buktikan \corref{intsubcor}.
\end{exercise}

\begin{exercise} \label{exercise:easyabsint}
Andaikan $f \colon [a,b] \to \R$ terbatas dan
memiliki sejumlah hingga diskontinuitas.
Tunjukkan bahwa, sebagai fungsi dari $x$, ekspresi $\babs{f(x)}$ terbatas dengan sejumlah hingga
diskontinuitas sehingga terintegralkan secara Riemann.  Kemudian tunjukkan 
\begin{equation*}
\abs{\int_a^b f(x)\,dx} \leq \int_a^b \babs{f(x)}\,dx .
\end{equation*}
\end{exercise}

\begin{exercise}[Menantang]
Tunjukkan bahwa fungsi
Thomae\index{fungsi Thomae} atau \myindex{fungsi popcorn}
(lihat \exampleref{popcornfunction:example})
terintegralkan secara Riemann.  Oleh karena itu, terdapat
fungsi yang diskontinu di semua bilangan rasional (suatu himpunan padat)
namun terintegralkan secara Riemann.

Secara lebih tepat,
definisikan $f \colon [0,1] \to \R$ dengan $f(0)=1$ dan, untuk $x > 0$, dengan
\begin{equation*}
f(x) \coloneqq 
\begin{cases}
\nicefrac{1}{k} & \text{jika } x=\nicefrac{m}{k} \text{ dengan }
m,k \in \N \text{ serta } m \text{ dan } k \text{ tidak memiliki pembagi bersama selain 1,} \\
0               & \text{jika } x \text{ irasional.}
\end{cases}
\end{equation*}
Tunjukkan bahwa $\int_0^1 f = 0$.
\end{exercise}


\begin{exnote}
Jika $I \subset \R$ suatu interval terbatas, maka
fungsi
\begin{equation*}
\varphi_I(x) \coloneqq
\begin{cases}
1 & \text{jika } x \in I, \\
0 & \text{jika tidak,}
\end{cases}
\end{equation*}
disebut \emph{\myindex{fungsi tangga elementer}}.
\end{exnote}

\begin{exercise} \label{exercise:stepfunctionintegrable}
Misalkan $I$ interval terbatas sebarang (anda harus mempertimbangkan semua jenis
interval: tertutup, terbuka, setengah terbuka) dan $a < b$.  Dengan
hanya menggunakan definisi integral,
tunjukkan bahwa
fungsi tangga elementer $\varphi_I$ terintegralkan
pada $[a,b]$, dan tentukan integralnya dalam bentuk $a$, $b$, serta
titik-titik ujung $I$.
\end{exercise}

\begin{exnote}
Suatu fungsi $f$ disebut
\emph{\myindex{fungsi tangga}} jika
dapat ditulis sebagai
\begin{equation*}
f(x) = \sum_{k=1}^n \alpha_k \varphi_{I_k} (x)
\end{equation*}
untuk beberapa bilangan real $\alpha_1,\alpha_2, \ldots, \alpha_n$
dan beberapa interval terbatas $I_1,I_2,\ldots,I_n$.
\end{exnote}

\begin{exercise}
Dengan menggunakan \exerciseref{exercise:stepfunctionintegrable}, tunjukkan bahwa suatu fungsi
tangga (lihat di atas) terintegralkan
pada setiap interval $[a,b]$.  Selanjutnya, tentukan integralnya dalam bentuk
$a$, $b$, titik-titik ujung $I_k$, dan $\alpha_k$.
\end{exercise}

\begin{exercise}
\label{exercise:boundedvariationintegrable}
Misalkan $f \colon [a,b] \to \R$ suatu fungsi.
\begin{enumerate}[a)]
\item
Tunjukkan bahwa jika $f$ naik, maka fungsi tersebut terintegralkan secara
Riemann.  Petunjuk: Gunakan partisi seragam; setiap subinterval memiliki panjang yang sama.
\item
Gunakan bagian a) untuk menunjukkan bahwa jika $f$ turun, maka fungsi tersebut terintegralkan secara
Riemann.
\item
Andaikan\footnote{$h$ semacam ini dikatakan memiliki \emph{\myindex{variasi terbatas}}.}
$h = f-g$, dengan $f$ dan $g$ fungsi-fungsi naik
pada $[a,b]$.  Tunjukkan bahwa $h$ terintegralkan secara Riemann.
\end{enumerate}
\end{exercise}

\begin{exercise}[Menantang] \label{exercise:hardabsint}
Andaikan $f \in \sR\bigl([a,b]\bigr)$.  Buktikan bahwa fungsi yang memetakan $x$ ke
$\babs{f(x)}$ juga terintegralkan secara Riemann pada $[a,b]$.
Kemudian tunjukkan pertidaksamaan yang sama seperti dalam \exerciseref{exercise:easyabsint}.
\end{exercise}

\begin{exercise}  \label{exercise:upperlowerlinineq}
Andaikan $f \colon [a,b] \to \R$ dan $g \colon [a,b] \to \R$
terbatas.
\begin{enumerate}[a)]
\item
Tunjukkan
$\overline{\int_a^b} (f+g) \leq \overline{\int_a^b}f+\overline{\int_a^b}g$ dan
$\underline{\int_a^b} (f+g) \geq
\underline{\int_a^b}f+\underline{\int_a^b}g$.
\item
Temukan contoh $f$ dan $g$ yang membuat
pertidaksamaan tersebut ketat.  Petunjuk: $f$ dan $g$ hendaknya tidak terintegralkan secara
Riemann.
\end{enumerate}
\end{exercise}

\begin{exercise}
Andaikan $f \colon [a,b] \to \R$ kontinu dan $g \colon \R \to \R$
kontinu Lipschitz.  Definisikan
\begin{equation*}
h(x) \coloneqq \int_a^b g(t-x) f(t) \, dt .
\end{equation*}
Buktikan bahwa $h$ kontinu Lipschitz.
\end{exercise}

\begin{exercise} \label{exercise:rieleblem}
Buktikan suatu versi dari \emph{\myindex{lemma Riemann--Lebesgue}}
(salah satu dari beberapa hasil yang menyandang nama tersebut):
Andaikan $f \colon [a,b] \to \R$ kontinu dan definisikan barisan
$\{ x_n \}_{n=1}^\infty$ dengan
\begin{equation*}
x_n \coloneqq \int_a^b f(t) \sin(nt) \, dt .
\end{equation*}
Buktikan bahwa $\lim\limits_{n\to\infty} x_n = 0$.
\end{exercise}

%%%%%%%%%%%%%%%%%%%%%%%%%%%%%%%%%%%%%%%%%%%%%%%%%%%%%%%%%%%%%%%%%%%%%%%%%%%%%%

\sectionnewpage
\section{Teorema dasar kalkulus}
\label{sec:ftc}

%mbxINTROSUBSECTION

\sectionnotes{1,5 kuliah}

Pada bagian ini kita membahas dan membuktikan
\emph{\myindex{teorema dasar kalkulus}}.
Seluruh kalkulus integral dibangun di atas teorema ini,
dan itulah sebabnya namanya demikian.
Teorema ini mengaitkan konsep integral dan turunan yang sepintas tidak berkaitan.
Teorema ini memberi tahu kita cara menghitung antiturunan suatu fungsi
menggunakan integral dan sebaliknya.

\subsection{Bentuk pertama teorema}

\begin{thm} \label{thm:FTCv1}
Misalkan $F \colon [a,b] \to \R$ merupakan fungsi kontinu yang diferensiabel
pada $(a,b)$.  Misalkan $f \in \sR\bigl([a,b]\bigr)$ sedemikian sehingga $f(x) = F'(x)$ untuk $x \in
(a,b)$.  Maka
\begin{equation*}
\int_a^b f = F(b)-F(a) .
\end{equation*}
\end{thm}

Teorema tersebut tidak sulit digeneralisasi agar mengizinkan berhingga banyak titik
di $[a,b]$ tempat $F$ tidak diferensiabel, asalkan fungsi tersebut tetap kontinu.
Generalisasi ini diserahkan sebagai latihan.

\begin{proof}
Misalkan $P = \{ x_0, x_1, \ldots, x_n \}$ merupakan partisi dari $[a,b]$.
Untuk setiap interval $[x_{i-1},x_i]$, gunakan
\hyperref[thm:mvt]{teorema nilai rata-rata} untuk memperoleh
$c_i \in (x_{i-1},x_i)$ sedemikian sehingga
\begin{equation*}
f(c_i) \Delta x_i = F'(c_i) (x_i - x_{i-1}) = F(x_i) - F(x_{i-1}) .
\end{equation*}
Lihat \figureref{fig:fundthmfig}, dan
perhatikan bahwa luas persegi panjang
ke-$i$ adalah
$F(x_{i})-F(x_{i-1})$,
sedangkan luas total ketiga persegi panjang yang digambarkan adalah
$F(x_{i+1})-F(x_{i-2})$.
Gagasannya adalah bahwa dengan mengambil subinterval yang semakin kecil,
luas total semua persegi panjang ini konvergen ke integral $f$.
\begin{myfigureht}
\myincludegraphics{fundthmfig}{%
Diagram grafik suatu fungsi yang digambar tebal gelap
dan diberi tanda y sama dengan f dari x sama dengan F aksen dari x,
beserta tiga subinterval koordinat x.
Subinterval tengah diberi label membentang dari
x subskrip kuantitas i dikurangi 1 hingga x subskrip i dan panjangnya Delta x subskrip i.
Titik c subskrip i ditandai di dalam subinterval ini dan sebuah
garis putus-putus ditarik vertikal ke atas hingga f dari c subskrip i, tempat garis itu berpotongan dengan
grafik f.  Sebuah persegi panjang berarsir dengan tinggi ini dan subinterval
sebagai alas digambar dan diberi label 'luas sama dengan f dari c subskrip i dikali Delta
x subskrip i sama dengan F dari x subskrip i dikurangi F dari x subskrip kuantitas i dikurangi 1'.
Dua subinterval lainnya serupa, kecuali pada yang kiri
kita mengganti i dengan i dikurangi 1 dan pada yang kanan mengganti i dengan
i ditambah 1.}
\caption{Teorema nilai rata-rata pada subinterval-subinterval suatu partisi yang
menghampiri luas di bawah kurva.\label{fig:fundthmfig}}
\end{myfigureht}

Dengan notasi dari definisi integral,
$m_i \leq f(c_i) \leq M_i$, dan dengan mengalikannya dengan $\Delta x_i$ diperoleh
\begin{equation*}
m_i \Delta x_i \leq F(x_i) - F(x_{i-1}) \leq M_i \Delta x_i .
\end{equation*}
Menjumlahkan terhadap $i = 1,2, \ldots, n$ menghasilkan
\begin{equation*}
\sum_{i=1}^n m_i \Delta x_i
\leq \sum_{i=1}^n \bigl(F(x_i) - F(x_{i-1}) \bigr)
\leq \sum_{i=1}^n M_i \Delta x_i .
\end{equation*}
Dalam jumlah di tengah, semua suku selain suku pertama dan terakhir saling menghapus,
sehingga tersisa $F(x_n)-F(x_0) = F(b)-F(a)$.  Jumlah di sebelah kiri
dan kanan masing-masing merupakan jumlah bawah dan jumlah atas.  Jadi
\begin{equation*}
L(P,f) \leq F(b)-F(a) \leq U(P,f) .
\end{equation*}
Kita mengambil supremum $L(P,f)$ atas semua partisi $P$, dan pertidaksamaan kiri
menghasilkan
\begin{equation*}
\underline{\int_a^b} f \leq F(b)-F(a) .
\end{equation*}
Serupa dengan itu, mengambil
infimum $U(P,f)$ atas semua partisi $P$ menghasilkan
\begin{equation*}
F(b)-F(a) \leq \overline{\int_a^b} f .
\end{equation*}
Karena $f$ terintegralkan Riemann, kita memperoleh
\begin{equation*}
\int_a^b f =
\underline{\int_a^b} f \leq F(b)-F(a) \leq \overline{\int_a^b} f
= \int_a^b f .
\end{equation*}
Semua pertidaksamaan tersebut harus berupa kesamaan, dan pembuktian selesai.
\end{proof}

Teorema tersebut digunakan untuk menghitung integral.  Andaikan kita mengetahui bahwa
fungsi $f(x)$ merupakan turunan dari suatu fungsi lain $F(x)$,
maka kita dapat menemukan bentuk eksplisit untuk $\int_a^b f$.

\begin{example}
Untuk menghitung
\begin{equation*}
\int_0^1 x^2 \,dx ,
\end{equation*}
kita perhatikan bahwa $x^2$ merupakan turunan dari $\frac{x^3}{3}$.
Teorema dasar mengatakan bahwa
\begin{equation*}
\int_0^1 x^2 \,dx =
\frac{1^3}{3}
-
\frac{0^3}{3}
= \frac{1}{3}.
\end{equation*}
\end{example}

\subsection{Bentuk kedua teorema}

Bentuk kedua teorema dasar memberi kita cara untuk menyelesaikan
persamaan diferensial $F'(x) = f(x)$, dengan $f$ merupakan fungsi yang diketahui
dan kita berusaha mencari $F$ yang memenuhi persamaan tersebut.

\begin{thm} \label{thm:FTCv2}
Misalkan $f \colon [a,b] \to \R$ merupakan fungsi terintegralkan Riemann.  Definisikan
\begin{equation*}
F(x) \coloneqq \int_a^x f .
\end{equation*}
Pertama, $F$ kontinu pada $[a,b]$.  Kedua,
jika $f$ kontinu di $c \in [a,b]$, maka $F$ diferensiabel di $c$
dan $F'(c) = f(c)$.
\end{thm}

\begin{proof}
Karena $f$ terbatas, terdapat $M > 0$
sedemikian sehingga $\babs{f(x)} \leq M$ untuk semua $x \in [a,b]$.  Andaikan $x,y \in [a,b]$
dengan $x > y$.  Maka
\begin{equation*}
\babs{F(x)-F(y)} =
\abs{\int_a^x f - \int_a^y f}
=
\abs{\int_y^x f}
\leq
M\sabs{x-y} .
\end{equation*}
Karena simetri, hal yang sama juga berlaku jika $x < y$.
Jadi $F$ kontinu secara Lipschitz dan karena itu kontinu.

Sekarang andaikan $f$ kontinu di $c$.
Misalkan diberikan $\epsilon > 0$.  Pilih $\delta > 0$ sedemikian sehingga,
untuk $x \in [a,b]$,
$\sabs{x-c} < \delta$ mengakibatkan $\babs{f(x)-f(c)} < \epsilon$.
Secara khusus,
untuk $x$ semacam itu, kita mempunyai
\begin{equation*}
f(c)-\epsilon < f(x) < f(c) + \epsilon.
\end{equation*}
Jadi jika $x > c$, maka
\begin{equation*}
\bigl(f(c)-\epsilon\bigr) (x-c) \leq \int_c^x f \leq
\bigl(f(c) + \epsilon\bigr)(x-c).
\end{equation*}
Jika $c > x$, arah pertidaksamaan tersebut berbalik.  Oleh karena itu,
dengan mengandaikan $x \neq c$, kita memperoleh
\begin{equation*}
f(c)-\epsilon
\leq
\frac{\int_c^{x} f}{x-c}
\leq
f(c)+\epsilon .
\end{equation*}
Karena
\begin{equation*}
\frac{F(x)-F(c)}{x-c}
=
\frac{\int_a^{x} f - \int_a^{c} f}{x-c}
=
\frac{\int_c^{x} f}{x-c} ,
\end{equation*}
kita mempunyai
\begin{equation*}
\abs{\frac{F(x)-F(c)}{x-c} - f(c)} \leq \epsilon .
\end{equation*}
Hasilnya pun mengikuti.  Pembaca dipersilakan memeriksa mengapa tidak menjadi masalah bahwa
pertidaksamaan yang kita peroleh hanya tidak ketat.
\end{proof}

Tentu saja, jika $f$ kontinu pada $[a,b]$, maka fungsi tersebut otomatis terintegralkan
Riemann, $F$ diferensiabel pada seluruh $[a,b]$, dan $F'(x) = f(x)$ untuk
semua $x \in [a,b]$.

\begin{remark} \label{remark:fundthmbase}
Bentuk kedua teorema dasar kalkulus tetap berlaku jika
kita mengambil $d \in [a,b]$ dan mendefinisikan
\begin{equation*}
F(x) \coloneqq \int_d^x f .
\end{equation*}
Artinya, kita dapat menggunakan sembarang titik di $[a,b]$ sebagai titik pangkal.  Pembuktiannya
diserahkan sebagai latihan.
\end{remark}

Mari kita lihat akibat suatu diskontinuitas sederhana.  Ambil $f(x) \coloneqq -1$ jika $x
< 0$, dan $f(x) \coloneqq 1$ jika $x \geq 0$.  Misalkan $F(x) \coloneqq \int_0^x f$.  Tidak
sulit untuk melihat bahwa $F(x) = \sabs{x}$.  Perhatikan bahwa $f$ diskontinu di
$0$ dan $F$ tidak diferensiabel di $0$.  Namun, kebalikan dari
teorema tersebut tidak
berlaku.
Misalkan $g(x) \coloneqq 0$ jika $x \neq 0$, dan $g(0) \coloneqq 1$.
Dengan mengambil $G(x) \coloneqq \int_0^x g$,
kita memperoleh $G(x) = 0$ untuk semua $x$.  Jadi $g$ diskontinu
di $0$, tetapi $G'(0)$ ada dan sama dengan 0.

Kesalahpahaman umum tentang integral di kalangan mahasiswa kalkulus adalah
menganggap integral yang hasilnya tidak dapat diberikan dalam bentuk tertutup sebagai sesuatu yang
cacat.  Hal ini tidak benar.  Kebanyakan integral yang kita tuliskan tidak dapat
dihitung dalam bentuk tertutup.  Bahkan beberapa integral yang kita anggap
berbentuk tertutup sebenarnya tidak demikian.  Kita mendefinisikan
logaritma alami sebagai antiturunan dari $\nicefrac{1}{x}$
yang memenuhi $\ln 1 = 0$:
\begin{equation*}
\ln x \coloneqq \int_1^x \frac{1}{s}\,ds .
\end{equation*}
Bagaimana komputer mencari nilai $\ln x$?
Salah satu caranya ialah menghampiri integral ini secara numerik.
Pada hakikatnya,
kita tidak benar-benar \myquote{menyederhanakan} $\int_1^x \frac{1}{s}\,ds$ dengan
menuliskan $\ln x$.  Kita sekadar memberi integral tersebut sebuah nama.
Jika kita memerlukan jawaban numerik,
mungkin pada akhirnya kita tetap melakukan
perhitungan dengan menghampiri suatu integral.
Pada bagian berikutnya, kita bahkan mendefinisikan fungsi eksponensial menggunakan
logaritma, yang kita definisikan melalui integral.

Fungsi umum lain yang didefinisikan oleh suatu integral yang tidak dapat
dievaluasi secara simbolik dalam fungsi-fungsi elementer
adalah fungsi $\operatorname{erf}$, yang didefinisikan sebagai
\begin{equation*}
\operatorname{erf}(x) \coloneqq \frac{2}{\sqrt{\pi}} \int_0^x e^{-s^2} \,ds .
\end{equation*}
Fungsi ini sering muncul dalam matematika terapan.  Fungsi tersebut semata-mata merupakan
antiturunan dari $\left(\nicefrac{2}{\sqrt{\pi}}\right) e^{-x^2}$
yang bernilai nol di nol.
Bentuk kedua teorema dasar memberi tahu kita bahwa fungsi tersebut dapat ditulis
sebagai integral.  Jika kita ingin menghitung suatu nilai tertentu, kita
menghampiri integral itu secara numerik.

\subsection{Penggantian variabel}

Salah satu teorema yang sering digunakan dalam kalkulus untuk menghitung integral adalah teorema
penggantian variabel; Anda mungkin mengenalnya sebagai
\emph{substitusi-$u$}\index{substitusi-u@substitusi-$u$}.
Agar mudah dihubungkan dengan materi kalkulus, kita juga akan
menamai variabel baru tersebut $u$, yaitu kita melakukan substitusi $u=g(x)$.
Ingat bahwa suatu fungsi diferensiabel secara kontinu jika
fungsi tersebut diferensiabel dan turunannya kontinu.

\begin{thm}[Penggantian variabel]
\index{teorema penggantian variabel}
Misalkan $g \colon [a,b] \to \R$ fungsi yang diferensiabel secara kontinu,
misalkan $f \colon [c,d] \to \R$ kontinu, dan andaikan
$g\bigl([a,b]\bigr) \subset [c,d]$.  Maka
\begin{equation*}
\int_a^b f\bigl(g(x)\bigr)\, g'(x)\, dx =
\int_{g(a)}^{g(b)} f(u)\, du .
\end{equation*}
\end{thm}

\begin{proof}
Karena $g$, $g'$, dan $f$ kontinu, $f\bigl(g(x)\bigr)\,g'(x)$
merupakan fungsi kontinu pada $[a,b]$, sehingga terintegralkan secara Riemann.
Demikian pula, $f$ terintegralkan pada setiap subinterval dari $[c,d]$.

Definisikan $F \colon [c,d] \to \R$ dengan
\begin{equation*}
F(y) \coloneqq \int_{g(a)}^{y} f(u)\,du .
\end{equation*}
Berdasarkan bentuk kedua teorema dasar
kalkulus
(lihat \remarkref{remark:fundthmbase} dan \exerciseref{secondftc:exercise}),
$F$ merupakan fungsi diferensiabel dan $F'(y) = f(y)$.  Terapkan aturan
rantai,
\begin{equation*}
\bigl( F \circ g \bigr)' (x) =
F'\bigl(g(x)\bigr) g'(x)
=
f\bigl(g(x)\bigr) g'(x) .
\end{equation*}
Perhatikan bahwa $F\bigl(g(a)\bigr) = 0$, lalu
gunakan bentuk pertama teorema dasar
untuk memperoleh
\begin{multline*}
%mbxSTARTIGNORE
\qquad %to center things more
%mbxENDIGNORE
\int_{g(a)}^{g(b)} f(u)\,du = F\bigl(g(b)\bigr) = F\bigl(g(b)\bigr)-F\bigl(g(a)\bigr)
\\
=
\int_a^b 
\bigl( F \circ g \bigr)' (x) \,dx
=
\int_a^b 
f\bigl(g(x)\bigr) g'(x)
\,dx .
%mbxSTARTIGNORE
\qquad %to center things more
%mbxENDIGNORE
\qedhere
\end{multline*}
\end{proof}

Teorema penggantian variabel sering digunakan untuk menghitung integral dengan mengubahnya
menjadi integral yang telah kita ketahui atau dapat kita hitung menggunakan teorema dasar
kalkulus.

\begin{example}
Turunan $\sin(x)$ adalah $\cos(x)$.
Dengan menggunakan $g(x) \coloneqq x^2$, kita hitung
\begin{equation*}
\int_0^{\sqrt{\pi}} x \cos(x^2) \, dx = \int_0^\pi \frac{\cos(u)}{2} \, du
=
\frac{1}{2}
\int_0^\pi \cos(u) \, du
=
\frac{
\sin(\pi) - \sin(0)
}{2}
=
0 .
\end{equation*}
\end{example}

Namun, ingatlah bahwa kita harus memenuhi hipotesis teorema tersebut.  Contoh
berikut menunjukkan mengapa kita tidak boleh sekadar
memindah-mindahkan simbol tanpa berpikir.
Kita harus memastikan bahwa simbol-simbol tersebut benar-benar bermakna.

\begin{example}
Perhatikan
\begin{equation*}
\int_{-1}^{1} \frac{\ln \sabs{x}}{x} \,dx .
\end{equation*}
Mungkin kita tergoda untuk mengambil $g(x) \coloneqq \ln \sabs{x}$.  Hitung $g'(x) =
\nicefrac{1}{x}$ dan coba tuliskan
\begin{equation*}
\int_{g(-1)}^{g(1)} u \,du = 
\int_{0}^{0} u \,du = 0. 
\end{equation*}
\myquote{Penyelesaian} ini salah dan tidak menunjukkan
bahwa kita dapat menghitung integral yang diberikan.  Masalah pertama adalah
$\frac{\ln \sabs{x}}{x}$ tidak kontinu pada $[-1,1]$.
Fungsi tersebut tidak terdefinisi di 0 dan tidak dapat dijadikan kontinu dengan menetapkan suatu nilai di
0.
Kedua, $\frac{\ln \sabs{x}}{x}$ bahkan tidak terintegralkan secara Riemann pada $[-1,1]$
(fungsi tersebut tak terbatas).
Integral yang kita tuliskan sama sekali tidak bermakna.
Terakhir, $g$ tidak kontinu
pada $[-1,1]$, apalagi diferensiabel secara kontinu.
\end{example}

\subsection{Latihan}

\begin{exercise}
Hitung
$\displaystyle
\frac{d}{dx} \biggl( \int_{-x}^x e^{s^2}\,ds \biggr)$.
\end{exercise}

\begin{exercise}
Hitung
$\displaystyle
\frac{d}{dx} \biggl( \int_{0}^{x^2} \sin(s^2)\,ds \biggr)$.
\end{exercise}

\begin{exercise}
Andaikan $F \colon [a,b] \to \R$ kontinu dan diferensiabel
pada $[a,b] \setminus S$, dengan $S$ suatu himpunan hingga.  Andaikan
terdapat $f \in \sR\bigl([a,b]\bigr)$ sedemikian sehingga $f(x) = F'(x)$ untuk $x \in [a,b]
\setminus S$.  Tunjukkan bahwa
$\int_a^b f = F(b)-F(a)$.
\end{exercise}

\begin{exercise} \label{secondftc:exercise}
Misalkan $f \colon [a,b] \to \R$ fungsi kontinu.  Pilih sebarang $c \in [a,b]$.
Definisikan
\begin{equation*}
F(x) \coloneqq \int_c^x f .
\end{equation*}
Buktikan bahwa $F$ diferensiabel dan $F'(x) = f(x)$ untuk semua $x \in
[a,b]$.
\end{exercise}

\begin{exercise}
Buktikan \emph{\myindex{integrasi parsial}}.  Artinya, andaikan $F$ dan
$G$ merupakan fungsi yang diferensiabel secara kontinu pada $[a,b]$.  Buktikan
\begin{equation*}
\int_a^b F(x)G'(x)\,dx
=
F(b)G(b)-F(a)G(a)
-
\int_a^b F'(x)G(x)\,dx .
\end{equation*}
\end{exercise}

\begin{exercise}
Andaikan $F$ dan $G$
diferensiabel secara kontinu\footnote{
Bandingkan hipotesis ini dengan \exerciseref{exercise:samediffconst}.}
dan didefinisikan
pada $[a,b]$
sedemikian sehingga $F'(x) = G'(x)$ untuk semua $x \in [a,b]$.
Dengan menggunakan teorema dasar kalkulus,
tunjukkan bahwa selisih $F$ dan $G$ merupakan konstanta.  Artinya, tunjukkan bahwa
terdapat $C \in \R$ sedemikian sehingga
$F(x)-G(x) = C$.
\end{exercise}

\begin{exnote}
Latihan berikut menunjukkan cara menggunakan integral untuk \myquote{memperhalus}
fungsi yang tidak diferensiabel.
\end{exnote}

\begin{exercise} \label{exercise:smoothingout}
Misalkan $f \colon [a,b] \to \R$ fungsi kontinu.  Misalkan $\epsilon > 0$
konstanta sedemikian sehingga $a+\epsilon < b-\epsilon$.  Untuk $x \in [a+\epsilon,b-\epsilon]$, definisikan
\begin{equation*}
g(x) \coloneqq \frac{1}{2\epsilon} \int_{x-\epsilon}^{x+\epsilon} f .
\end{equation*}
\begin{enumerate}[a)]
\item
Tunjukkan bahwa $g$ diferensiabel dan tentukan turunannya.
\item
Misalkan $f$ diferensiabel dan tetapkan $x \in (a,b)$ (ambil $\epsilon$
cukup kecil).  Apa yang terjadi pada $g'(x)$ ketika $\epsilon$ semakin kecil?
\item
Tentukan $g$ untuk $f(x) \coloneqq \sabs{x}$, $\epsilon = 1$ (Anda dapat mengandaikan
$[a,b]$ cukup besar).
\end{enumerate}
\end{exercise}

\begin{exercise}
Andaikan $f \colon [a,b] \to \R$ kontinu dan
$\int_a^x f = \int_x^b f$ untuk semua $x \in [a,b]$.  Tunjukkan bahwa $f(x) = 0$
untuk semua $x \in [a,b]$.
\end{exercise}

\begin{exercise}
Andaikan $f \colon [a,b] \to \R$ kontinu dan
$\int_a^x f = 0$ untuk setiap $x$ rasional dalam $[a,b]$.  Tunjukkan bahwa $f(x) = 0$
untuk semua $x \in [a,b]$.
\end{exercise}

\begin{samepage}
\begin{exercise}
Fungsi $f$ disebut \emph{\myindex{fungsi ganjil}} jika $f(x) = -f(-x)$,
dan $f$ disebut \emph{\myindex{fungsi genap}} jika $f(x) = f(-x)$.  Misalkan $a >
0$.  Andaikan $f$ kontinu.  Buktikan:
\begin{enumerate}[a)]
\item
Jika $f$ ganjil, maka $\int_{-a}^a f = 0$.
\item
Jika $f$ genap, maka $\int_{-a}^a f = 2 \int_0^a f$.
\end{enumerate}
\end{exercise}
\end{samepage}

\begin{exercise}
\leavevmode
\begin{enumerate}[a)]
\item
Tunjukkan bahwa $f(x) \coloneqq \sin(\nicefrac{1}{x})$
terintegralkan pada setiap interval (Anda dapat mendefinisikan $f(0)$ sebagai nilai apa pun).
\item
Hitung $\int_{-1}^1 \sin(\nicefrac{1}{x})\,dx$ (perhatikan
diskontinuitasnya).
\end{enumerate}
\end{exercise}

\begin{exercise}[menggunakan \sectionref{sec:monotonefunc}]
\leavevmode
\begin{enumerate}[a)]
\item
Andaikan $f \colon [a,b] \to \R$ monoton naik.
Berdasarkan \propref{prop:monotoneintegrable},
%\exerciseref{exercise:boundedvariationintegrable},
$f$ terintegralkan secara Riemann.  Tunjukkan bahwa jika $f$ diskontinu di $c \in
(a,b)$, maka $F(x) \coloneqq \int_a^x f$ tidak diferensiabel di $c$.
\item
Dalam \exerciseref{exercise:increasingfuncdiscatQ}, Anda mengonstruksi fungsi monoton naik
$f \colon [0,1] \to \R$ yang diskontinu di setiap
$x \in [0,1] \cap \Q$.  Gunakan $f$ ini untuk mengonstruksi fungsi $F(x)$ yang
kontinu pada $[0,1]$; untuk setiap $x \in [0,1] \cap \Q$, fungsi tersebut tidak diferensiabel di titik itu.
\end{enumerate}
\end{exercise}

\begin{exercise}
Untuk sebarang $\ell \in \N$,
tunjukkan bahwa limit berikut ada dan tentukan nilainya:
\begin{equation*}
\lim_{n\to\infty} \sum_{k=1}^n \frac{k^\ell}{n^{\ell+1}}
\end{equation*}
\end{exercise}

%%%%%%%%%%%%%%%%%%%%%%%%%%%%%%%%%%%%%%%%%%%%%%%%%%%%%%%%%%%%%%%%%%%%%%%%%%%%%%

\sectionnewpage
\section{Logaritma dan eksponensial}
\label{sec:logandexp}

%mbxINTROSUBSECTION

\sectionnotes{1 kuliah (opsional, memerlukan bagian opsional
\sectionref{sec:limitatinf},
\sectionref{sec:monotonefunc},
\sectionref{sec:ift})}

Kita kini memiliki perangkat yang diperlukan untuk mendefinisikan dengan tepat fungsi eksponensial
dan
logaritma yang telah Anda kenal dari kalkulus.  Kita mulai dengan
perpangkatan.
Jika $n$ bilangan bulat positif, kita definisikan
\begin{equation*}
x^n \coloneqq \underbrace{x \cdot x \cdots x}_{n \text{ kali}} .
\end{equation*}
Wajar untuk mendefinisikan $x^0 \coloneqq 1$.
Untuk bilangan bulat negatif, tetapkan $x^{-n} \coloneqq \nicefrac{1}{x^n}$ selama
$x \neq 0$.
Selanjutnya, andaikan $x > 0$.
Definisikan
$x^{1/n}$ sebagai
satu-satunya akar pangkat $n$ yang positif.  Terakhir, untuk bilangan
rasional $\nicefrac{n}{m}$ (dalam bentuk paling sederhana), definisikan
\begin{equation*}
x^{n/m} \coloneqq {\bigl(x^{1/m}\bigr)}^n .
\end{equation*}
Tidak sulit untuk menunjukkan bahwa
kita memperoleh bilangan yang sama apa pun
representasi $\nicefrac{n}{m}$ yang digunakan, sehingga sebenarnya kita tidak perlu menggunakan
bentuk paling sederhana.

Namun, apa yang dimaksud dengan $\sqrt{2}^{\sqrt{2}}$?  Atau,
$x^y$ secara umum?  Secara khusus, bagaimana $e^x$ didefinisikan untuk setiap~$x$?
Dan bagaimana kita menyelesaikan $y=e^x$ untuk~$x$?
Bagian ini menjawab pertanyaan-pertanyaan tersebut dan beberapa pertanyaan lain.

\subsection{Logaritma}
\index{logaritma}

Lebih mudah jika kita mendefinisikan logaritma terlebih dahulu.
Mari kita tunjukkan bahwa
terdapat tepat satu fungsi dengan sifat-sifat yang sesuai; baru setelah itu
kita akan menyebutnya \emph{logaritma}.

\begin{prop}
Terdapat tepat satu fungsi $L \colon (0,\infty) \to \R$ sedemikian sehingga
\begin{enumerate}[(i)]
\item \label{it:log:i}
$L(1) = 0$.
\item \label{it:log:ii}
$L$ diferensiabel dan $L'(x) = \nicefrac{1}{x}$.
\item \label{it:log:iii}
$L$ naik tegas, bijektif, dan
\begin{equation*}
\lim_{x\to 0} L(x) = -\infty , \qquad \text{dan} \qquad
\lim_{x\to \infty} L(x) = \infty .
\end{equation*}
\item \label{it:log:iv}
$L(xy) = L(x)+L(y)$ untuk semua $x,y \in (0,\infty)$.
\item \label{it:log:v}
Jika $q$ bilangan rasional dan $x > 0$, maka
$L(x^q) = q L(x)$.
\end{enumerate}
\end{prop}

\begin{proof}
Untuk membuktikan keberadaan, kita definisikan suatu kandidat dan tunjukkan bahwa kandidat tersebut memenuhi
semua sifat.  Tetapkan
\begin{equation*}
L(x) \coloneqq \int_1^x \frac{1}{t}\,dt .
\end{equation*}

Jelas bahwa \ref{it:log:i} berlaku.  Sifat \ref{it:log:ii} berlaku
berdasarkan bentuk kedua teorema dasar kalkulus
(\thmref{thm:FTCv2}).

Untuk membuktikan sifat \ref{it:log:iv},
kita mengganti variabel dengan $u=yt$ dan memperoleh
\begin{equation*}
L(x) =
\int_1^{x} \frac{1}{t}\,dt
=
\int_y^{xy} \frac{1}{u}\,du
=
\int_1^{xy} \frac{1}{u}\,du
-
\int_1^{y} \frac{1}{u}\,du
=
L(xy)-L(y) .
\end{equation*}

Mari kita buktikan \ref{it:log:iii}.
Sifat \ref{it:log:ii}, bersama dengan fakta bahwa $L'(x) = \nicefrac{1}{x} > 0$
untuk $x > 0$, menyiratkan bahwa $L$
naik tegas dan karenanya injektif.
Mari kita tunjukkan bahwa $L$ surjektif.
Karena $\nicefrac{1}{t} \geq \nicefrac{1}{2}$ ketika $t \in [1,2]$,
\begin{equation*}
L(2) = \int_1^2 \frac{1}{t} \,dt \geq \nicefrac{1}{2} .
\end{equation*}
Berdasarkan induksi, \ref{it:log:iv} menyiratkan bahwa untuk $n \in \N$,
\begin{equation*}
L(2^n) = L(2) + L(2) + \cdots + L(2) = n L(2) .
\end{equation*}
Diberikan $y > 0$,
berdasarkan \hyperref[thm:arch:i]{sifat Archimedes} bilangan real
(perhatikan bahwa $L(2) > 0$), terdapat $n \in \N$ sedemikian sehingga
$L(2^n) > y$.  Kemudian,
\hyperref[IVT:thm]{teorema nilai antara}
memberikan $x_1 \in (1,2^n)$ sedemikian sehingga $L(x_1) = y$.  Jadi,
semua bilangan dalam $(0,\infty)$ berada dalam citra $L$.
Karena $L$ naik, $L(x) > y$ untuk semua $x > 2^n$, sehingga
\begin{equation*}
\lim_{x\to\infty} L(x) = \infty .
\end{equation*}
Selanjutnya,
$0 = L(\nicefrac{x}{x}) = L(x) + L(\nicefrac{1}{x})$, sehingga
$L(x) = - L(\nicefrac{1}{x})$.  Dengan menggunakan $x=2^{-n}$, seperti di atas kita memperoleh
bahwa $L$ mencapai semua bilangan negatif.  Selain itu,
\begin{equation*}
\lim_{x \to 0} L(x) = 
\lim_{x \to 0} -L(\nicefrac{1}{x})
=
\lim_{x \to \infty} -L(x)
=  - \infty .
\end{equation*}
Dalam limit-limit tersebut, perhatikan bahwa hanya $x > 0$ yang berada dalam domain $L$.

Mari kita buktikan \ref{it:log:v}.
Tetapkan $x > 0$.
Seperti di atas, \ref{it:log:iv} menyiratkan
$L(x^n) = n L(x)$
untuk semua $n \in \N$.
Kita telah memperoleh bahwa
$L(x) = - L(\nicefrac{1}{x})$,
sehingga $L(x^{-n}) = - L(x^n) = -n L(x)$.  Selanjutnya, untuk $m \in \N$,
\begin{equation*}
L(x) = L\Bigl({(x^{1/m})}^m\Bigr) = m L\bigl(x^{1/m}\bigr) .
\end{equation*}
Dengan menggabungkan semuanya untuk $n \in \Z$ dan $m \in \N$, kita memperoleh
$L(x^{n/m}) = n L(x^{1/m}) = (\nicefrac{n}{m}) L(x)$.

Ketunggalan mengikuti dari sifat \ref{it:log:i} dan
\ref{it:log:ii}.  Berdasarkan bentuk pertama
teorema dasar kalkulus (\thmref{thm:FTCv1}),
\begin{equation*}
L(x) = \int_1^x \frac{1}{t}\,dt
\end{equation*}
merupakan satu-satunya fungsi sedemikian sehingga $L(1) = 0$ dan $L'(x) = \nicefrac{1}{x}$.
\end{proof}

Setelah membuktikan
bahwa terdapat tepat satu fungsi dengan sifat-sifat tersebut,
kita mendefinisikan \emph{\myindex{logaritma}}, yang terkadang disebut
\emph{\myindex{logaritma alami}}:
\glsadd{not:ln}
\begin{equation*}
\ln(x) \coloneqq L(x) .
\end{equation*}
Lihat \figureref{fig:log}.
Matematikawan biasanya menulis $\log(x)$, sedangkan mahasiswa kalkulus lebih akrab
dengan $\ln(x)$.  Untuk semua keperluan praktis, hanya ada satu logaritma:
logaritma alami.  Lihat \exerciseref{exercise:otherlogbases}.
\begin{myfigureht}
\myincludegraphics{logfig}{%
Grafik y sama dengan ln dari x digambar dengan garis tebal yang berawal curam
dengan nilai y negatif di sebelah kiri x sama dengan 1,
memotong sumbu x di 1, dan semakin landai
ke arah kanan.
Grafik y sama dengan 1 per x digambar dengan garis putus-putus; garis ini
selalu positif, berawal tinggi di sebelah kiri, dan
semakin landai ke arah kanan.
Daerah antara garis putus-putus dan sumbu x
untuk x antara 1 dan 4 diarsir, dan nilai grafik
ln pada x sama dengan 4 ditandai sebagai 'luas daerah yang diarsir sama dengan ln dari 4'.}
\caption{Grafik $\ln(x)$ bersama $\nicefrac{1}{x}$, yang menunjukkan nilai
$\ln(4)$.\label{fig:log}}
\end{myfigureht}

\subsection{Eksponensial}
\index{eksponensial}

Seperti halnya logaritma, kita mendefinisikan eksponensial melalui sebuah daftar
sifat.

\begin{prop}
Terdapat tepat satu fungsi $E \colon \R \to (0,\infty)$ sedemikian sehingga
\begin{enumerate}[(i)]
\item \label{it:exp:i}
$E(0) = 1$.
\item \label{it:exp:ii}
$E$ diferensiabel dan $E'(x) = E(x)$.
\item \label{it:exp:iii}
$E$ naik tegas, bijektif, dan
\begin{equation*}
\lim_{x\to -\infty} E(x) = 0 \qquad \text{dan} \qquad
\lim_{x\to \infty} E(x) = \infty .
\end{equation*}
\item \label{it:exp:iv}
$E(x+y) = E(x)E(y)$ untuk semua $x,y \in \R$.
\item \label{it:exp:v}
Jika $q \in \Q$ dan $x \in \R$, maka
$E(qx) = {E(x)}^q$.
\end{enumerate}
\end{prop}

\begin{proof}
Sekali lagi, kita membuktikan keberadaan fungsi semacam itu dengan mendefinisikan suatu kandidat
dan membuktikan bahwa kandidat tersebut memenuhi semua sifat.
Fungsi $L = \ln$ yang didefinisikan di atas memiliki invers.  Misalkan $E$ merupakan
fungsi invers dari $L$.  Sifat \ref{it:exp:i} berlaku langsung.

Sifat \ref{it:exp:ii} mengikuti
dari teorema fungsi invers, khususnya
dari \lemmaref{lemma:ift}:  $L$~memenuhi
semua hipotesis lema tersebut, sehingga
\begin{equation*}
E'(x) = \frac{1}{L'\bigl(E(x)\bigr)} = E(x) .
\end{equation*}

Mari kita tinjau sifat \ref{it:exp:iii}.
Fungsi $E$ naik tegas karena
$E'(x) = E(x) > 0$.  Karena $E$ merupakan fungsi invers dari $L$, fungsi tersebut juga pasti
bijektif.
Untuk menentukan limit-limitnya, kita gunakan fakta bahwa
$E$ naik tegas dan surjektif ke $(0,\infty)$.
Untuk setiap $M > 0$, terdapat $x_0$ sedemikian sehingga
$E(x_0) = M$ dan $E(x) \geq M$ untuk semua $x \geq x_0$.
Demikian pula, untuk setiap $\epsilon > 0$, terdapat
$x_0$ sedemikian sehingga $E(x_0) = \epsilon$ dan
$E(x) < \epsilon$ untuk semua $x < x_0$.
Oleh karena itu,
\begin{equation*}
\lim_{x\to -\infty} E(x) = 0 \qquad \text{dan} \qquad
\lim_{x\to \infty} E(x) = \infty .
\end{equation*}

Untuk membuktikan sifat \ref{it:exp:iv}, kita gunakan sifat yang bersesuaian
untuk logaritma.
Ambil $x, y \in \R$.
Karena $L$ bijektif, terdapat $a$ dan $b$ sedemikian sehingga $x = L(a)$ dan $y = L(b)$.  Maka
\begin{equation*}
E(x+y) =
E\bigl(L(a)+L(b)\bigr) = 
E\bigl(L(ab)\bigr) = ab = E(x)E(y)  .
\end{equation*}

Sifat \ref{it:exp:v} juga mengikuti dari sifat $L$ yang bersesuaian.
Diberikan $x \in \R$, misalkan $a$ sedemikian sehingga $x = L(a)$.  Maka
\begin{equation*}
E(qx) = E\bigl(qL(a)\bigr)
=
E\bigl(L(a^q)\bigr) = a^q = {E(x)}^q .
\end{equation*}

Ketunggalan mengikuti dari
\ref{it:exp:i} dan
\ref{it:exp:ii}.
Misalkan $E$ dan $F$
merupakan dua fungsi yang memenuhi
\ref{it:exp:i} dan \ref{it:exp:ii}.
\begin{equation*}
\frac{d}{dx} \Bigl( F(x)E(-x) \Bigr)
=
F'(x)E(-x) - E'(-x)F(x)
=
F(x)E(-x) - E(-x)F(x) = 0 .
\end{equation*}
Oleh karena itu, berdasarkan \propref{prop:derzeroconst},
$F(x)E(-x) = F(0)E(-0) = 1$ untuk semua $x \in \R$.
Dengan melakukan perhitungan tersebut untuk $F = E$,
kita memperoleh $E(x)E(-x) = 1$.
Maka
\begin{equation*}
0 = 1-1 = F(x)E(-x) - E(x)E(-x) = \bigl(F(x)-E(x)\bigr) E(-x) .
\end{equation*}
%Since $E(x)E(-x) = 1$,
Terakhir, $E(-x) \neq 0$\footnote{%
Bagaimanapun, kodomain $E$ adalah $(0,\infty)$.
Namun, $E(-x) \neq 0$ juga mengikuti
dari $E(x)E(-x) = 1$.  Oleh karena itu, kita dapat membuktikan ketunggalan $E$
dengan hanya mengandaikan \ref{it:exp:i} dan \ref{it:exp:ii}, bahkan untuk fungsi $E \colon \R
\to \R$.}
untuk semua $x \in \R$.
Jadi,
$F(x)-E(x) = 0$ untuk semua $x$, dan pembuktian selesai.
\end{proof}

Setelah membuktikan bahwa $E$ tunggal, kita mendefinisikan fungsi
\emph{\myindex{eksponensial}} (lihat \figureref{fig:exp}) sebagai
\glsadd{not:exp}
\begin{equation*}
\exp(x) \coloneqq E(x) .
\end{equation*}
\begin{myfigureht}
\myincludegraphics{expfig}{%
Medan kemiringan digambar pada bidang xy untuk x yang kira-kira berada antara minus 1 dan 1
serta y antara 0 dan 4.  Medan kemiringan tersebut berupa kisi ruas-ruas garis pendek dengan
kemiringan tertentu.  Dalam gambar ini, kemiringannya sama pada setiap baris.
Semua ruas mengarah ke atas (positif) dan
semakin curam ketika kita bergerak ke atas.  Grafik fungsi
eksponensial juga ditampilkan.  Grafik tersebut menunjukkan fungsi naik yang berawal dari sebelah kiri
dan mengikuti kemiringan-kemiringan yang diberikan:
Berawal di atas sumbu x, grafik melewati titik (0,1), lalu menjadi
semakin curam ketika keluar dari gambar di sebelah kanan.}
\caption{Grafik $e^x$, bersama medan kemiringan dengan kemiringan $y$ di
setiap titik $(x,y)$.
Persamaan $\frac{d}{dx} e^x = e^x$ berarti bahwa $y=e^x$ mengikuti kemiringan-kemiringan tersebut.\label{fig:exp}}
\end{myfigureht}

Jika $y \in \Q$ dan $x > 0$, maka
\begin{equation*}
x^y = \exp\bigl(\ln(x^y)\bigr) = \exp\bigl(y\ln(x)\bigr) .
\end{equation*}
Sekarang kita dapat memberi makna pada perpangkatan $x^y$ untuk sembarang $y \in \R$;
jika $x > 0$ dan $y$ irasional, definisikan
\glsadd{not:pow}
\begin{equation*}
x^y \coloneqq \exp\bigl(y\ln(x)\bigr) .
\end{equation*}
Karena $\exp$ kontinu, $x^y$ merupakan fungsi kontinu terhadap $y$.
Oleh karena itu, kita akan
memperoleh hasil yang sama jika mengambil suatu barisan bilangan rasional
$\{ y_n \}_{n=1}^\infty$
yang konvergen ke $y$ dan mendefinisikan $x^y = \lim_{n\to\infty} x^{y_n}$.

Definisikan bilangan $e$,
yang disebut \emph{\myindex{bilangan Euler}} atau
\emph{\myindex{basis logaritma alami}}, sebagai
\glsadd{not:e}
\begin{equation*}
e \coloneqq \exp(1) .
\end{equation*}
Mari kita benarkan notasi $e^x$ untuk $\exp(x)$:
\begin{equation*}
e^x = \exp\bigl(x \ln(e) \bigr) = \exp(x) .
\end{equation*}

Sifat-sifat logaritma dan eksponensial juga berlaku untuk
pangkat irasional.  Buktinya langsung.

\begin{prop}
Misalkan $x, y \in \R$.
\begin{enumerate}[(i)]
\item
$\exp(xy) = {\bigl(\exp(x)\bigr)}^y$.
\item
Jika $x > 0$, maka $\ln(x^y) = y \ln (x)$.
\end{enumerate}
\end{prop}

\begin{remark}
Ada cara-cara ekuivalen lain untuk mendefinisikan eksponensial dan logaritma.
Salah satu cara yang umum adalah mendefinisikan $E$ sebagai solusi persamaan diferensial
$E'(x) = E(x)$, $E(0) = 1$.  Lihat \exampleref{example:picardexponential}
untuk sketsa pendekatan tersebut.  Pendekatan lainnya
adalah mendefinisikan fungsi eksponensial melalui
deret pangkat; lihat \exampleref{example:exponentialbypowerseries}.
\end{remark}

\begin{remark}
Kita telah membuktikan ketunggalan fungsi $L$ dan $E$ hanya dari
sifat-sifat $L(1)=0$, $L'(x) = \nicefrac{1}{x}$
serta syarat yang ekuivalen untuk eksponensial,
$E'(x) = E(x)$, $E(0) = 1$.  Keberadaan juga mengikuti
hanya dari sifat-sifat tersebut.
Sebagai alternatif, ketunggalan juga mengikuti
dari hukum-hukum eksponen; lihat latihan.
\end{remark}

\subsection{Latihan}

\begin{exercise}
Untuk bilangan real $y$ dan $b > 0$, definisikan $f \colon (0,\infty) \to \R$
dan $g \colon \R \to \R$ dengan $f(x) \coloneqq x^y$ dan $g(x) \coloneqq b^x$.  Tunjukkan bahwa $f$
dan $g$ diferensiabel dan tentukan turunan keduanya.
\end{exercise}

\begin{samepage}
\begin{exercise} \label{exercise:otherlogbases}
Misalkan $b > 0$, $b\neq 1$.
\begin{enumerate}[a)]
\item
Tunjukkan bahwa untuk setiap $y > 0$, terdapat tepat satu bilangan $x$
sedemikian sehingga $y = b^x$.  Definisikan
\emph{\myindex{logaritma dengan basis $b$}},
$\log_b \colon (0,\infty) \to \R$, dengan
$\log_b(y) \coloneqq x$.
\item
Tunjukkan bahwa $\log_b(x) = \frac{\ln(x)}{\ln(b)}$.
\item
Buktikan bahwa jika $c > 0$, $c \neq 1$, maka
$\log_b(x) = \frac{\log_c(x)}{\log_c(b)}$.
\item
Buktikan bahwa $\log_b(xy) =
\log_b(x)+\log_b(y)$ dan $\log_b(x^y) = y \log_b(x)$.
\end{enumerate}
\end{exercise}
\end{samepage}

\begin{exercise}[memerlukan \sectionref{sec:taylor}]
Gunakan \hyperref[thm:taylor]{teorema Taylor} untuk menelaah suku sisa dan menunjukkan bahwa untuk
semua $x \in \R$
\begin{equation*}
e^x = \sum_{n=0}^\infty \frac{x^n}{n!} .
\end{equation*}
Petunjuk: Jangan mendiferensialkan deret tersebut suku demi suku (kecuali Anda membuktikan bahwa
hal itu dapat dilakukan).
\end{exercise}

\begin{exercise}
Gunakan rumus jumlah deret geometri untuk menunjukkan (untuk $t \neq -1$)
\begin{equation*}
1-t+t^2-\cdots+{(-1)}^n t^n = \frac{1}{1+t} - \frac{{(-1)}^{n+1}t^{n+1}}{1+t}.
\end{equation*}
Dengan menggunakan fakta ini, tunjukkan bahwa
\begin{equation*}
\ln (1+x) = \sum_{n=1}^\infty \frac{{(-1)}^{n+1}x^n}{n} 
\end{equation*}
untuk semua $x \in (-1,1]$ (perhatikan bahwa $x=1$ tercakup).  Terakhir,
tentukan limit deret harmonik berselang-seling
\begin{equation*}
\sum_{n=1}^\infty \frac{{(-1)}^{n+1}}{n} = 1 - \frac{1}{2} +
\frac{1}{3} - \frac{1}{4} + \cdots
\end{equation*}
\end{exercise}

\begin{exercise}
Tunjukkan bahwa
\begin{equation*}
e^x = \lim_{n\to\infty} {\left( 1 + \frac{x}{n} \right)}^n .
\end{equation*}
Petunjuk: Ambil logaritmanya.\\
Catatan: Bentuk
${\left( 1 + \frac{x}{n} \right)}^n$ muncul dalam perhitungan bunga
majemuk.  Bentuk ini menyatakan jumlah uang dalam rekening bank setelah 1 tahun
jika pada awalnya 1 dolar disimpan dengan tingkat bunga $x$ (misalnya, $x=0.01$ untuk
$1\%$) dan bunganya dimajemukkan $n$
kali selama tahun tersebut.  Eksponensial $e^x$ merupakan hasil pemajemukan
kontinu.
\end{exercise}

\begin{samepage}
\begin{exercise}
\leavevmode
\begin{enumerate}[a)]
\item
Buktikan bahwa untuk semua bilangan bulat $n \geq 2$,
\begin{equation*}
\sum_{k=2}^{n}
\frac{1}{k}
\leq
\ln (n)
\leq
\sum_{k=1}^{n-1}
\frac{1}{k} .
\end{equation*}
\item
Buktikan bahwa limit
\begin{equation*}
\gamma \coloneqq \lim_{n\to\infty}
\left( \left( \sum_{k=1}^{n} \frac{1}{k} \right) - \ln (n) \right)
\end{equation*}
ada.  Konstanta $\gamma$ ini dikenal sebagai
\emph{\myindex{konstanta Euler--Mascheroni}}%
\footnote{Dinamai menurut matematikawan Swiss
\href{https://en.wikipedia.org/wiki/Leonhard_Euler}{Leonhard Euler}
(1707--1783)
dan matematikawan Italia
\href{https://en.wikipedia.org/wiki/Lorenzo_Mascheroni}{Lorenzo Mascheroni}
(1750--1800).}.  Belum diketahui apakah $\gamma$ rasional atau tidak.
Nilainya kira-kira $\gamma \approx 0.5772$.
\end{enumerate}
\end{exercise}
\end{samepage}

\begin{exercise}
Tunjukkan bahwa
\begin{equation*}
\lim_{x\to\infty} \frac{\ln(x)}{x} = 0 .
\end{equation*}
\end{exercise}

\begin{exercise}
Tunjukkan bahwa $e^x$ \emph{\myindex{konveks}}; dengan kata lain, tunjukkan bahwa
jika $a < b$ dan $a \leq x \leq b$, maka
$e^x \leq e^a \frac{b-x}{b-a} + e^b \frac{x-a}{b-a}$.
\end{exercise}

\begin{exercise}
Dengan menggunakan logaritma, tentukan
\begin{equation*}
%\lim_{n\to\infty} {\left( 1 + \nicefrac{1}{n} \right)}^n = e .
\lim_{n\to\infty} n^{1/n} .
\end{equation*}
\end{exercise}

\begin{exercise}
Tunjukkan bahwa $E(x) = e^x$ merupakan satu-satunya fungsi kontinu sedemikian sehingga
$E(x+y) = E(x)E(y)$ dan $E(1) = e$.   Dengan cara serupa, buktikan bahwa $L(x) = \ln(x)$
merupakan satu-satunya fungsi kontinu
yang terdefinisi untuk $x$ positif sedemikian sehingga $L(xy) = L(x)+L(y)$
dan $L(e) = 1$.
\end{exercise}

\begin{exercise}[memerlukan \sectionref{sec:taylor}]\label{exercise:nonanalytic}
Karena $(e^x)' = e^x$, mudah dilihat bahwa $e^x$
\myindex{diferensiabel tak hingga kali}\index{diferensiabel!tak hingga kali}
(memiliki turunan dari setiap orde).  Definisikan fungsi $f \colon \R \to \R$.
\begin{equation*}
f(x) \coloneqq \begin{cases}
e^{-1/x} & \text{jika } x > 0, \\
0 & \text{jika } x \leq 0.
\end{cases}
\end{equation*}
\begin{enumerate}[a)]
\item
Buktikan bahwa untuk setiap $m \in \N$,
\begin{equation*}
\lim_{x \to 0^+} \frac{e^{-1/x}}{x^m} = 0 .
\end{equation*}
\item
Buktikan bahwa $f$ diferensiabel tak hingga kali.
\item
Hitung deret Taylor untuk $f$ di titik asal, yaitu
\begin{equation*}
\sum_{k=0}^\infty
\frac{f^{(k)}(0)}{k!}x^k .
\end{equation*}
Tunjukkan bahwa deret tersebut konvergen, tetapi tidak konvergen ke $f(x)$
untuk setiap $x > 0$.
\end{enumerate}
\end{exercise}

%%%%%%%%%%%%%%%%%%%%%%%%%%%%%%%%%%%%%%%%%%%%%%%%%%%%%%%%%%%%%%%%%%%%%%%%%%%%%%

\sectionnewpage
\section{Integral tak wajar}
\label{sec:impropriemann}

%mbxINTROSUBSECTION

\sectionnotes{2--3 pertemuan kuliah (bagian opsional, dapat dilewati tanpa masalah, 
memerlukan \sectionref{sec:limitatinf} yang opsional)}

Sering kali kita perlu mengintegralkan pada seluruh
garis bilangan real, atau pada interval tak terbatas berbentuk $[a,\infty)$ atau
$(-\infty,b]$.  Kita mungkin juga ingin mengintegralkan fungsi tak terbatas
yang terdefinisi pada interval terbuka terbatas $(a,b)$.
Untuk interval atau fungsi semacam itu, integral Riemann tidak terdefinisi,
tetapi kita tetap akan menuliskan
integralnya dengan mengikuti gagasan dalam \lemmaref{lemma:boundedimpriemann}.
Integral-integral ini disebut \emph{\myindex{integral tak wajar}}
dan merupakan limit
dari integral, bukan integral itu sendiri.

\begin{defn}
Andaikan $f \colon [a,b) \to \R$ adalah fungsi (tidak harus terbatas)
yang terintegralkan secara Riemann pada $[a,c]$ untuk setiap $a < c < b$.  Kita definisikan
\begin{equation*}
\int_a^b f \coloneqq \lim_{c \to b^-} \int_a^{c} f
\end{equation*}
jika limit tersebut ada.

Andaikan $f \colon [a,\infty) \to \R$ adalah fungsi sedemikian sehingga
$f$ terintegralkan secara Riemann pada $[a,c]$ untuk setiap $a < c < \infty$.
Kita definisikan
\begin{equation*}
\int_a^\infty f \coloneqq \lim_{c \to \infty} \int_a^c f
\end{equation*}
jika limit tersebut ada.

Jika limit tersebut ada, kita katakan bahwa integral tak wajar itu
\emph{konvergen}\index{konvergen!integral tak wajar}.
Jika limit tersebut tidak ada, kita katakan bahwa integral tak wajar itu
\emph{divergen}\index{divergen!integral tak wajar}.

Dengan cara serupa, kita definisikan integral tak wajar untuk titik ujung kiri.
Definisi ini diserahkan kepada pembaca.
\end{defn}

Untuk titik ujung berhingga $b$, jika $f$ terbatas, maka
\lemmaref{lemma:boundedimpriemann} mengatakan bahwa kita tidak mendefinisikan sesuatu yang baru.  Hal yang baru adalah
kita dapat menerapkan definisi ini pada fungsi tak terbatas.
Kumpulan contoh berikut
begitu berguna sehingga kita menyatakannya sebagai suatu proposisi.

\begin{prop}[Uji-$p$ untuk integral]%
\index{uji-p untuk integral@uji-$p$ untuk integral}
\label{impropriemann:ptest}
Integral tak wajar
\begin{equation*}
\int_1^\infty \frac{1}{x^p} \,dx
\end{equation*}
konvergen ke $\frac{1}{p-1}$ jika $p > 1$ dan divergen jika $p \leq 1$.

Integral tak wajar
\begin{equation*}
\int_0^1 \frac{1}{x^p} \,dx
\end{equation*}
konvergen ke $\frac{1}{1-p}$ jika $p < 1$ dan divergen jika $p \geq 1$.
\end{prop}

\begin{proof}
Bukti diperoleh dengan menerapkan
\hyperref[thm:FTCv1]{teorema dasar kalkulus}.
Mari kita buktikan kasus $p > 1$ untuk titik ujung kanan tak hingga dan
menyerahkan sisanya kepada pembaca.  Petunjuk: Kasus $p=1$ harus ditangani
secara terpisah.

Andaikan $p > 1$.  Dengan menggunakan
teorema dasar tersebut, kita peroleh
\begin{equation*}
\int_1^b \frac{1}{x^p} \,dx
=
\int_1^b x^{-p} \,dx
=
\frac{b^{-p+1}}{-p+1}
-
\frac{1^{-p+1}}{-p+1}
=
\frac{-1}{(p-1)b^{p-1}}
+
\frac{1}{p-1} .
\end{equation*}
Karena $p > 1$, kita mempunyai $p-1 > 0$.  Dengan mengambil limit saat $b \to \infty$,
kita peroleh bahwa $\frac{1}{b^{p-1}}$ menuju 0.  Hasil yang diinginkan pun terbukti.
\end{proof}

Kita menyatakan proposisi berikut tentang \myquote{ekor} hanya untuk satu jenis
integral tak wajar, meskipun buktinya langsung
dan sama untuk jenis integral tak wajar lainnya.

\begin{prop} \label{impropriemann:tail}
Misalkan $f \colon [a,\infty) \to \R$ adalah fungsi
yang terintegralkan secara Riemann pada $[a,b]$ untuk setiap $b > a$.
Untuk setiap $b > a$, integral
$\int_b^\infty f$ konvergen jika dan hanya jika $\int_a^\infty f$
konvergen; dalam hal ini,
\begin{equation*}
\int_a^\infty f
=
\int_a^b f +
\int_b^\infty f .
\end{equation*}
\end{prop}

\begin{proof}
Misalkan $c > b$.  Maka
\begin{equation*}
\int_a^c f
=
\int_a^b f +
\int_b^c f .
\end{equation*}
Mengambil limit $c \to \infty$ menyelesaikan bukti.
\end{proof}

Fungsi tak negatif lebih mudah ditangani,
seperti ditunjukkan oleh proposisi berikut.
Latihan-latihan akan menunjukkan bahwa proposisi ini
hanya berlaku untuk fungsi tak negatif.
Terdapat proposisi analog
untuk semua jenis integral tak wajar lainnya, dan pembuktiannya diserahkan kepada
mahasiswa.

\begin{prop} \label{impropriemann:possimp}
Andaikan $f \colon [a,\infty) \to \R$ tak negatif ($f(x)
\geq 0$ untuk setiap $x$) dan
$f$ terintegralkan secara Riemann pada $[a,b]$ untuk setiap $b > a$.
\begin{enumerate}[(i)]
\item 
\begin{equation*}
\int_a^\infty f = \sup \left\{ \int_a^x f : x \geq a \right\} .
\end{equation*}
\item
Andaikan $\{ x_n \}_{n=1}^\infty$ adalah barisan dalam $[a,\infty)$
dengan $\lim_{n\to\infty} x_n = \infty$.  Maka
$\int_a^\infty f$ konvergen jika dan hanya jika $\lim_{n\to\infty} \int_a^{x_n} f$ ada; dalam
hal ini,
\begin{equation*}
\int_a^\infty f = \lim_{n\to\infty} \int_a^{x_n} f .
\end{equation*}
\end{enumerate}
\end{prop}

Dalam butir pertama, kita mengizinkan nilai $\infty$ pada
supremum untuk menunjukkan bahwa integral divergen menuju tak hingga.

\begin{proof}
Kita mulai dengan butir pertama.
Karena $f$ tak negatif,
$\int_a^x f$ meningkat sebagai fungsi dari $x$.
Jika supremumnya tak hingga, maka untuk setiap $M \in \R$
terdapat $N$ sedemikian sehingga $\int_a^N f \geq M$.  Karena $\int_a^x f$
meningkat, $\int_a^x f \geq M$ untuk setiap $x \geq N$.  Jadi,
$\int_a^\infty f$ divergen menuju tak hingga.

Selanjutnya, andaikan supremumnya berhingga, katakanlah
$A \coloneqq \sup \left\{ \int_a^x f : x \geq a \right\}$.
Untuk setiap $\epsilon > 0$, terdapat $N$ sedemikian sehingga
$A - \int_a^N f < \epsilon$.  Karena $\int_a^x f$ meningkat,
maka
$A - \int_a^x f < \epsilon$ untuk setiap $x \geq N$, sehingga
$\int_a^\infty f$ konvergen ke $A$.

Sekarang mari kita tinjau butir kedua.
Jika $\int_a^\infty f$ konvergen, maka setiap barisan $\{ x_n \}_{n=1}^\infty$ yang menuju
tak hingga memenuhi kesimpulan.  Bagian yang memerlukan kecermatan adalah
membuktikan arah sebaliknya.  Andaikan $\{ x_n \}_{n=1}^\infty$ sedemikian sehingga
$\lim_{n\to\infty} x_n = \infty$, dan limit berikut ada serta bernilai $A$:
\begin{equation*}
\lim_{n\to\infty} \int_a^{x_n} f = A
\end{equation*}
Diberikan $\epsilon > 0$, pilih $N$ sedemikian sehingga untuk
setiap $n \geq N$, berlaku
$A - \epsilon < \int_a^{x_n} f < A + \epsilon$.
Karena $\int_a^x f$ meningkat sebagai fungsi dari $x$, untuk setiap
$x \geq x_N$
\begin{equation*}
A - \epsilon < \int_a^{x_N} f \leq \int_a^x f .
\end{equation*}
Karena $\{ x_n \}_{n=1}^\infty$ menuju $\infty$, untuk setiap 
$x$ terdapat $x_m$ sedemikian sehingga $m \geq N$ dan $x \leq x_m$.  Maka
\begin{equation*}
\int_a^{x} f \leq \int_a^{x_m} f < A + \epsilon .
\end{equation*}
Khususnya, untuk setiap $x \geq x_N$, berlaku
$\abs{\int_a^{x} f - A} < \epsilon$.
\end{proof}

\begin{prop}[Uji perbandingan untuk integral tak wajar]%
\index{uji perbandingan untuk integral tak wajar}
Misalkan
$f \colon [a,\infty) \to \R$ dan
$g \colon [a,\infty) \to \R$ adalah fungsi
yang terintegralkan secara Riemann pada $[a,b]$ untuk setiap $b > a$.   Andaikan
bahwa untuk setiap $x \geq a$,
\begin{equation*}
\babs{f(x)} \leq g(x) .
\end{equation*}
\begin{enumerate}[(i)]
\item Jika $\int_a^\infty g$ konvergen, maka $\int_a^\infty f$ konvergen,
dan dalam hal ini 
$\abs{\int_a^\infty f} \leq \int_a^\infty g$.
\item Jika $\int_a^\infty f$ divergen, maka $\int_a^\infty g$ divergen.
\end{enumerate}
\end{prop}

\begin{proof}
Kita mulai dengan butir pertama.
Untuk setiap $b$ dan $c$ sedemikian sehingga $a \leq b \leq c$, berlaku 
$-g(x) \leq f(x) \leq g(x)$, sehingga
\begin{equation*}
\int_b^c -g \leq \int_b^c f \leq \int_b^c g  .
\end{equation*}
Dengan kata lain, $\abs{\int_b^c f} \leq \int_b^c g$.

Diberikan $\epsilon > 0$.  Berdasarkan
\propref{impropriemann:tail},
\begin{equation*}
\int_a^\infty g =
\int_a^b g +
\int_b^\infty g .
\end{equation*}
Karena $\int_a^b g$ menuju
$\int_a^\infty g$ saat $b$ menuju tak hingga,
$\int_b^\infty g$ menuju 0 saat $b$ menuju tak hingga.  Pilih $B$
sedemikian sehingga
\begin{equation*}
\int_B^\infty g < \epsilon .
\end{equation*}
Karena $g$ tak negatif, jika $B \leq b < c$, maka
$\int_b^c g < \epsilon$ juga.
Misalkan $\{ x_n \}_{n=1}^\infty$ adalah barisan dalam $[a,\infty)$ yang menuju tak hingga.  Ambil $M$ sedemikian sehingga
$x_n \geq B$ untuk setiap $n \geq M$.  Ambil $n, m \geq M$
dengan $x_n \leq x_m$,
\begin{equation*}
\abs{\int_a^{x_m} f - \int_a^{x_n} f} 
=
\abs{\int_{x_n}^{x_m} f} 
\leq \int_{x_n}^{x_m} g < \epsilon .
\end{equation*}
Jadi, barisan $\bigl\{ \int_a^{x_n} f \bigr\}_{n=1}^\infty$ adalah barisan Cauchy dan karena itu konvergen.

Kita perlu menunjukkan bahwa limit tersebut tunggal.  Andaikan $\{ x_n \}_{n=1}^\infty$ adalah barisan
yang menuju tak hingga sedemikian sehingga
$\bigl\{ \int_a^{x_n} f \bigr\}_{n=1}^\infty$ konvergen ke $L_1$, dan $\{
y_n \}_{n=1}^\infty$ adalah barisan
yang menuju tak hingga sedemikian sehingga
$\bigl\{ \int_a^{y_n} f \bigr\}_{n=1}^\infty$ konvergen ke $L_2$.  Maka terdapat $n$ sedemikian
sehingga
$\babs{\int_a^{x_n} f - L_1} < \epsilon$ dan 
$\babs{\int_a^{y_n} f - L_2} < \epsilon$.  Kita juga dapat mengandaikan $x_n \geq B$
dan $y_n \geq B$.  Setelah menukar label kedua barisan jika perlu, andaikan pula $x_n \leq y_n$.  Maka
\begin{equation*}
\sabs{L_1 - L_2} \leq
\abs{L_1 - \int_a^{x_n} f}
+
\abs{\int_a^{x_n} f- \int_a^{y_n} f}
+
\abs{\int_a^{y_n} f - L_2}
<
\epsilon
+
\abs{\int_{x_n}^{y_n} f}
+
\epsilon
<
3 \epsilon.
\end{equation*}
Karena $\epsilon > 0$ sebarang, $L_1 = L_2$, sehingga
$\int_a^\infty f$ konvergen.
Di atas telah kita tunjukkan bahwa $\abs{\int_a^c f} \leq \int_a^c g$ untuk setiap $c > a$.
Dengan mengambil limit $c \to \infty$, butir pertama terbukti.

Butir kedua tidak lain adalah kontrapositif dari butir pertama.
\end{proof}

\begin{example}
Integral tak wajar
\begin{equation*}
\int_0^\infty \frac{\sin(x^2)(x+2)}{x^3+1} \,dx
\end{equation*}
konvergen.

Bukti:  Perhatikan bahwa kita hanya perlu menunjukkan
bahwa integral dari 1 hingga tak hingga tersebut konvergen.
Untuk $x \geq 1$ kita peroleh
\begin{equation*}
\abs{\frac{\sin(x^2)(x+2)}{x^3+1}}
\leq
\frac{x+2}{x^3+1}
\leq \frac{x+2}{x^3} \leq
\frac{x+2x}{x^3} \leq \frac{3}{x^2} .
\end{equation*}
Kemudian
\begin{equation*}
\int_1^\infty \frac{3}{x^2}\,dx
=
3 \int_1^\infty \frac{1}{x^2}\,dx
%=
%\lim_{c\to\infty} \int_1^c \frac{3}{x^2} \,dx
=
3 .
\end{equation*}
Jadi, dengan uji perbandingan dan uji ekor, integral semula
konvergen.
\end{example}

\begin{example}
Kita harus berhati-hati ketika melakukan manipulasi formal pada integral
tak wajar.
Integral
\begin{equation*}
\int_2^\infty \frac{2}{x^2-1}\,dx
\end{equation*}
konvergen berdasarkan uji perbandingan, sekali lagi dengan menggunakan $\nicefrac{1}{x^2}$.  Namun, jika kita
tergoda untuk menulis
\begin{equation*}
\frac{2}{x^2-1} = 
\frac{1}{x-1}
-
\frac{1}{x+1} 
\end{equation*}
dan mencoba mengintegralkan setiap bagian secara terpisah, kita tidak akan berhasil.
\emph{Tidak} benar bahwa kita dapat memisahkan integral tak
wajar menjadi dua; kita tidak dapat memisahkan limitnya.
\begin{equation*}
\begin{split}
\int_2^\infty \frac{2}{x^2-1} \,dx &=
\lim_{b\to \infty} \int_2^b \frac{2}{x^2-1} \,dx
\\
&=
\lim_{b\to \infty}
\left(
\int_2^b \frac{1}{x-1}\,dx
-
\int_2^b \frac{1}{x+1}\,dx
\right)
\\
&\neq
\int_2^\infty \frac{1}{x-1}\,dx
-
\int_2^\infty \frac{1}{x+1}\,dx .
\end{split}
\end{equation*}
Baris terakhir dalam perhitungan tersebut bahkan tidak bermakna.  Kedua
integral itu divergen menuju tak hingga, sebab kita dapat
menerapkan uji perbandingan secara tepat dengan
$\nicefrac{1}{x}$.  Kita memperoleh $\infty - \infty$.
\end{example}

Sekarang andaikan kita perlu mengambil limit di kedua titik ujung.

\begin{defn}
Andaikan $f \colon (a,b) \to \R$ suatu fungsi
yang terintegralkan secara Riemann pada $[c,d]$ untuk semua $c$, $d$
sedemikian sehingga $a < c < d < b$, maka kita definisikan
\begin{equation*}
\int_a^b f \coloneqq \lim_{c \to a^+} \, \lim_{d \to b^-} \, \int_{c}^{d} f
\end{equation*}
jika limit-limit tersebut ada.

Andaikan $f \colon \R \to \R$ suatu fungsi sedemikian sehingga
$f$ terintegralkan secara Riemann pada semua interval terbatas $[a,b]$.  Maka
kita definisikan
\begin{equation*}
\int_{-\infty}^\infty f \coloneqq \lim_{c \to -\infty} \, \lim_{d \to \infty} \, \int_c^d f
\end{equation*}
jika limit-limit tersebut ada.

Dengan cara serupa, kita definisikan integral tak wajar dengan satu titik ujung di tak hingga
dan satu titik ujung berhingga yang tak wajar.  Definisi ini diserahkan kepada pembaca.
\end{defn}

Kita harus selalu berhati-hati dengan limit ganda.  Definisi
di atas menyatakan bahwa mula-mula kita mengambil limit ketika $d$ menuju $b$ atau
$\infty$ untuk $c$ yang tetap, lalu kita mengambil limit terhadap $c$.
Kita perlu membuktikan bahwa dalam hal ini tidak menjadi masalah limit mana yang
kita hitung terlebih dahulu.

\begin{example}
\begin{equation*}
\int_{-\infty}^\infty \frac{1}{1+x^2} \, dx
=
\lim_{a \to -\infty} \, \lim_{b \to \infty} \,
\int_{a}^b \frac{1}{1+x^2} \, dx
=
\lim_{a \to -\infty} \, \lim_{b \to \infty}
\bigl( \arctan(b) - \arctan(a) \bigr)
=
\pi .
\end{equation*}
\end{example}

Dalam definisi tersebut, urutan limit selalu dapat dipertukarkan jika limit-limit itu
ada.  Mari kita nyatakan dan buktikan fakta ini hanya untuk limit pada tak hingga.

\begin{prop}
Andaikan $f \colon \R \to \R$ terintegralkan pada setiap interval terbatas
$[a,b]$.
Maka 
\begin{equation*}
\lim_{a \to -\infty} \, \lim_{b \to \infty} \, \int_a^b f
\quad \text{konvergen} \qquad \text{jika dan hanya jika} \qquad
\lim_{b \to \infty}
\,
\lim_{a \to -\infty}
\,
\int_a^b f
\quad
\text{konvergen,}
\end{equation*}
dalam hal ini kedua
ekspresi tersebut sama.  Jika salah satu dari kedua
ekspresi itu konvergen, maka integral tak wajar tersebut konvergen dan
\begin{equation*}
\lim_{a\to\infty}
\int_{-a}^a f
=
\int_{-\infty}^\infty f .
\end{equation*}
\end{prop}

\begin{proof}
Tanpa mengurangi keumuman, andaikan $a < 0$ dan $b > 0$.  Andaikan
ekspresi pertama konvergen.  Maka
\begin{equation*}
\begin{split}
\lim_{a \to -\infty} \, \lim_{b \to \infty} \, \int_a^b f
& =
\lim_{a \to -\infty} \, \lim_{b \to \infty}
\left(
\int_a^0 f
+
\int_0^b f
\right)
=
\left(
\lim_{a \to -\infty}
\int_a^0 f
\right)
+
\left(
 \lim_{b \to \infty}
\int_0^b f
\right) \\
& = 
 \lim_{b \to \infty}
\left(
\left(
\lim_{a \to -\infty}
\int_a^0 f
\right) 
+
\int_0^b f
\right)
=
 \lim_{b \to \infty} \,
\lim_{a \to -\infty}
\left(
\int_a^0 f
+
\int_0^b f
\right)  .
\end{split}
\end{equation*}
Perhitungan serupa menunjukkan arah sebaliknya.  Oleh karena itu, jika
salah satu ekspresi konvergen, maka integral tak wajar tersebut konvergen
dan
\begin{multline*}
\int_{-\infty}^\infty f
=
\lim_{a \to -\infty} \, \lim_{b \to \infty} \, \int_a^b f
=
\left(
\lim_{a \to -\infty}
\int_a^0 f
\right)
+
\left(
 \lim_{b \to \infty}
\int_0^b f
\right)
\\
=
\left(
\lim_{a \to \infty}
\int_{-a}^0 f
\right)
+
\left(
 \lim_{a \to \infty}
\int_0^a f
\right)
=
\lim_{a \to \infty}
\left(
\int_{-a}^0 f
+
\int_0^a f
\right)
=
\lim_{a \to \infty}
\int_{-a}^a f .
\end{multline*}
\end{proof}

\begin{example}
Di sisi lain, Anda harus berhati-hati untuk
mengambil setiap limit secara terpisah sebelum mengetahui konvergensinya.  Misalkan
$f(x) = \frac{x}{\sabs{x}}$ untuk $x \neq 0$ dan $f(0) = 0$.
Jika $a < 0$ dan $b > 0$, maka
\begin{equation*}
\int_{a}^b f
=
\int_{a}^0 f
+
\int_{0}^b f
=
a+b .
\end{equation*}
Untuk setiap $a < 0$ yang tetap, limit ketika $b \to \infty$ adalah tak hingga.
Jadi, bahkan limit pertama pun tidak ada,
dan integral takwajar $\int_{-\infty}^\infty f$ tidak konvergen.
Di sisi lain, jika $a > 0$, maka
\begin{equation*}
\int_{-a}^{a} f
=
(-a)+a = 0 .
\end{equation*}
Oleh karena itu,
\begin{equation*}
\lim_{a\to\infty}
\int_{-a}^{a} f
= 0 .
\end{equation*}
\end{example}

\begin{example}
Salah satu contoh yang patut diingat untuk integral takwajar
adalah \emph{\myindex{fungsi sinc}}%
\footnote{Singkatan dari bahasa Latin: \emph{sinus cardinalis}}.
Fungsi ini cukup sering muncul
baik dalam matematika murni maupun matematika terapan.  Definisikan
\begin{equation*}
\operatorname{sinc}(x) \coloneqq
\begin{cases}
\frac{\sin(x)}{x} & \text{jika } x \neq 0 , \\
1 & \text{jika } x = 0 .
\end{cases}
\end{equation*}
\begin{myfigureht}
\myincludegraphics{sincfig}{%
Grafik fungsi sinc pada interval dari minus 4 pi hingga 4 pi.
Grafik berosilasi menyerupai fungsi sinus di sebelah kanan
titik asal, dan menyerupai negatif dari fungsi sinus di sebelah kiri
titik asal.  Akar fungsi di titik asal tidak lagi ada dan nilai fungsi tepat 1
di titik asal.  Osilasinya dengan cepat menjadi semakin kecil
ketika kita menjauhi titik asal.  Fungsi tersebut jelas genap.}
\caption{Fungsi sinc.\label{figsinc}}
\end{myfigureht}

Tidaklah sulit untuk menunjukkan bahwa
fungsi sinc kontinu di nol, tetapi hal itu
belum penting saat ini.  Hal yang penting adalah
\begin{equation*}
\int_{-\infty}^\infty \operatorname{sinc}(x) \,dx = \pi ,
\qquad \text{sedangkan} \qquad
\int_{-\infty}^\infty \babs{\operatorname{sinc}(x)} \,dx = \infty .
\end{equation*}
Integral fungsi sinc merupakan analog kontinu dari
deret harmonik berselang-seling $\sum_{n=1}^\infty \nicefrac{{(-1)}^n}{n}$, sedangkan integral
nilai mutlaknya menyerupai deret harmonik biasa $\sum_{n=1}^\infty \nicefrac{1}{n}$.
Secara khusus, konvergensi integral tersebut harus dibuktikan secara langsung,
bukan dengan menggunakan uji perbandingan.

Kita tidak akan membuktikan pernyataan pertama secara persis.  Mari kita buktikan saja
bahwa integral fungsi sinc konvergen, tanpa perlu mempersoalkan
nilai limitnya secara persis.  Karena $\frac{\sin(-x)}{-x} = \frac{\sin(x)}{x}$, cukup
ditunjukkan bahwa
\begin{equation*}
\int_{2\pi}^\infty \frac{\sin(x)}{x}\,dx
\end{equation*}
konvergen.  Dengan cara ini, kita
juga menghindari $x=0$ sehingga pembuktiannya menjadi lebih sederhana.

Untuk setiap $n \in \N$, jika $x \in [\pi 2n, \pi (2n+1)]$, kita peroleh
\begin{equation*}
\frac{\sin(x)}{\pi (2n+1)}
\leq
\frac{\sin(x)}{x}
\leq
\frac{\sin(x)}{\pi 2n} ,
\end{equation*}
karena $\sin(x) \geq 0$.  Untuk $x \in [\pi (2n+1), \pi (2n+2)]$,
\begin{equation*}
\frac{\sin(x)}{\pi (2n+1)}
\leq
\frac{\sin(x)}{x}
\leq
\frac{\sin(x)}{\pi (2n+2)} ,
\end{equation*}
karena $\sin(x) \leq 0$.

Berdasarkan teorema dasar kalkulus,
\begin{equation*}
%mbxlatex \begin{aligned}
\frac{2}{\pi (2n+1)}
=
\int_{\pi 2n}^{\pi (2n+1)}
\frac{\sin(x)}{\pi (2n+1)}
\,dx
%mbxlatex &
\leq
\int_{\pi 2n}^{\pi (2n+1)}
\frac{\sin(x)}{x}
\,dx
%mbxlatex \\
%mbxlatex &
\leq
\int_{\pi 2n}^{\pi (2n+1)}
\frac{\sin(x)}{\pi 2n}
\,dx
=
\frac{1}{\pi n} .
%mbxlatex \end{aligned}
\end{equation*}
Dengan cara serupa,
\begin{equation*}
\frac{-2}{\pi (2n+1)}
\leq
\int_{\pi (2n+1)}^{\pi (2n+2)}
\frac{\sin(x)}{x}
\,dx
\leq
\frac{-1}{\pi (n+1)} .
\end{equation*}
Dengan menjumlahkan keduanya, kita peroleh
\begin{equation*}
%mbxlatex \begin{aligned}
0
=
\frac{2}{\pi (2n+1)}
+
\frac{-2}{\pi (2n+1)}
%mbxlatex &
\leq
\int_{2\pi n}^{2\pi (n+1)}
\frac{\sin(x)}{x}
\,dx
%mbxlatex \\
%mbxlatex &
\leq
\frac{1}{\pi n} 
+
\frac{-1}{\pi (n+1)} 
=
\frac{1}{\pi n(n+1)} .
%mbxlatex \end{aligned}
\end{equation*}
Lihat \figureref{fig:sincbound}.
\begin{myfigureht}
\myincludegraphics{sincbound}{%
Ditampilkan grafik beberapa fungsi pada interval
dari pi kali 2 n hingga pi kali buka kurung 2 n tambah 2 tutup kurung.
Semuanya menyerupai gelombang sinus yang amplitudonya mengecil.
Garis tebal gelap adalah grafik fungsi sinus x dibagi x
pada seluruh interval.
Seperti semua fungsi lainnya, nilainya nol di titik ujung kiri,
titik ujung kanan, dan titik tengah pi kali buka kurung 2 n
tambah 1 tutup kurung.  Garis titik-titik pada subinterval pertama (separuh kiri),
yang sangat dekat dengan dan berada di bawah garis penuh, menyatakan
fungsi sinus x dibagi pi kali buka kurung 2 n tambah 1 tutup kurung.
Garis putus-putus pada subinterval pertama, yang sangat dekat dengan dan berada di atas
garis penuh, menyatakan fungsi sinus x dibagi pi kali 2 n.
Daerah antara garis putus-putus dan sumbu x pada subinterval pertama
diarsir.
Pada separuh kanan interval, garis titik-titik menyatakan
fungsi yang sama seperti sebelumnya, dan sekali lagi berada di bawah garis penuh.
Garis putus-titik yang sedikit berada di atas garis penuh
kini menyatakan fungsi sinus x dibagi pi kali buka kurung
2 n tambah 2 tutup kurung.  Daerah di atas garis putus-titik diarsir.}
\caption{Batas untuk $\int_{2\pi n}^{2\pi (n+1)} \frac{\sin(x)}{x} \,dx$ dengan menggunakan
integral yang diarsir (luas bertanda
$\frac{1}{\pi n} 
+
\frac{-1}{\pi (n+1)}$).\label{fig:sincbound}}
\end{myfigureht}

Untuk $k \in \N$, 
\begin{equation*}
\int_{2\pi}^{2k\pi} \frac{\sin(x)}{x} \,dx
=
\sum_{n=1}^{k-1}
\int_{2\pi n}^{2\pi (n+1)} \frac{\sin(x)}{x} \,dx 
\leq
\sum_{n=1}^{k-1}
\frac{1}{\pi n(n+1)} .
\end{equation*}
Kita memperoleh jumlah-jumlah parsial dari suatu deret dengan suku-suku positif.
Deret tersebut
konvergen karena
$\sum_{n=1}^\infty \frac{1}{\pi n (n+1)}$ merupakan deret konvergen.  Jadi,
sebagai suatu barisan,
\begin{equation*}
\lim_{k\to \infty} \int_{2\pi}^{2k\pi} \frac{\sin(x)}{x} \,dx
=L \leq
\sum_{n=1}^{\infty}
\frac{1}{\pi n(n+1)} < \infty .
\end{equation*}

Misalkan $M > 2\pi$ sembarang, dan misalkan $k \in \N$
merupakan bilangan bulat terbesar sedemikian sehingga $2k\pi \leq M$.
Untuk $x \in [2k\pi,M]$, kita memiliki
$\frac{-1}{2k\pi} \leq \frac{\sin(x)}{x} \leq \frac{1}{2k\pi}$, sehingga
\begin{equation*}
\abs{\int_{2k\pi}^{M} \frac{\sin(x)}{x} \,dx }  \leq
\frac{M-2k\pi}{2k\pi} \leq \frac{1}{k} .
\end{equation*}
Karena $k$ adalah nilai $k$ terbesar yang memenuhi $2k\pi \leq M$,
ketika $M\in \R$ menuju tak hingga, $k \in \N$ juga menuju tak hingga.

Selanjutnya,
\begin{equation*}
\int_{2\pi}^M \frac{\sin(x)}{x}\,dx
=
\int_{2\pi}^{2k\pi} \frac{\sin(x)}{x} \,dx
+
\int_{2k\pi}^{M} \frac{\sin(x)}{x} \,dx .
\end{equation*}
Ketika $M$ menuju tak hingga,
suku pertama pada
ruas kanan menuju $L$,
dan suku kedua pada
ruas kanan
menuju nol.  Jadi,
\begin{equation*}
\int_{2\pi}^\infty \frac{\sin(x)}{x} \,dx = L .
%\leq \sum_{n=1}^{\infty}
%\frac{1}{\pi n(n+1)} < \infty .
\end{equation*}

Integral dua sisi dari fungsi sinc juga ada, sebagaimana telah disebutkan di atas.
Pernyataan lainnya---bahwa integral
nilai mutlak fungsi sinc divergen---kita serahkan sebagai latihan.
\end{example}

\subsection{Uji integral untuk deret}

Teorema dasar
kalkulus dapat digunakan untuk membuktikan bahwa suatu deret konvergen dan untuk menaksir jumlahnya.

\begin{prop}[Uji integral]\index{uji integral untuk deret}
Andaikan $f \colon [k,\infty) \to \R$ merupakan fungsi tak negatif yang
menurun, dengan $k \in \Z$.  Maka
\begin{equation*}
\sum_{n=k}^\infty f(n)
\quad \text{konvergen} \qquad \text{jika dan hanya jika}
\qquad
\int_k^\infty f
\quad \text{konvergen}.
\end{equation*}
Dalam hal ini,
\begin{equation*}
\int_k^\infty f
\leq
\sum_{n=k}^\infty f(n)
\leq
f(k)+
\int_k^\infty f .
\end{equation*}
\end{prop}
Lihat \figureref{fig:integraltest} untuk ilustrasi dengan $k=1$.
Berdasarkan \propref{prop:monotoneintegrable},
$f$ terintegralkan
pada setiap interval $[k,b]$ untuk semua $b > k$, sehingga pernyataan proposisi tersebut
bermakna tanpa hipotesis keterintegralan tambahan.
\begin{myfigureht}
\myincludegraphics{integraltest}{%
Diagram persegi panjang Darboux untuk fungsi menurun,
dimulai dari x sama dengan 1 dan berlanjut dalam interval
sepanjang 1 hingga x sama dengan 10 di sebelah kanan.  Karena fungsi
tersebut menurun, grafik fungsi melalui
sudut kiri atas setiap persegi panjang utuh dan juga sudut
kanan atas setiap persegi panjang yang diarsir.}
\caption{Luas di bawah kurva,
$\int_1^\infty f$, dibatasi dari bawah
oleh luas persegi panjang yang diarsir,
$f(2)+f(3)+f(4)+\cdots$, dan dibatasi dari atas
oleh luas seluruh persegi panjang, baik bagian yang diarsir maupun yang tidak,
$f(1)+f(2)+f(3)+\cdots$.\label{fig:integraltest}}
\end{myfigureht}

\begin{proof}
Misalkan $\ell, m \in \Z$ sedemikian sehingga $m > \ell \geq k$.
Karena $f$ menurun, kita memiliki
$\int_{n}^{n+1} f \leq f(n) \leq \int_{n-1}^{n} f$.  Oleh karena itu,
\begin{equation} \label{impropriemann:eqseries}
\int_\ell^m f
=
\sum_{n=\ell}^{m-1} \int_{n}^{n+1} f
\leq
\sum_{n=\ell}^{m-1} f(n)
\leq
f(\ell) +
\sum_{n=\ell+1}^{m-1} \int_{n-1}^{n} f
\leq
f(\ell)+
\int_\ell^{m-1} f .
\end{equation}

Pertama, andaikan $\int_k^\infty f$ konvergen dan
diberikan $\epsilon > 0$.
Seperti sebelumnya, karena $f$ tak negatif, terdapat
$L \in \N$ sedemikian sehingga jika $\ell \geq L$, maka
$\int_\ell^{m} f < \nicefrac{\epsilon}{2}$ untuk semua $m \geq \ell$.
Fungsi
$f$ harus menurun menuju nol (mengapa?), jadi pilih $L$ cukup besar sehingga
untuk $\ell \geq L$, berlaku $f(\ell) < \nicefrac{\epsilon}{2}$.
Dengan demikian, untuk $m > \ell \geq L$, berdasarkan
\eqref{impropriemann:eqseries}, kita memiliki
\begin{equation*}
\sum_{n=\ell}^{m} f(n)
\leq
f(\ell)+
\int_\ell^{m} f < \nicefrac{\epsilon}{2} + \nicefrac{\epsilon}{2} = \epsilon .
\end{equation*}
Oleh karena itu, deret tersebut bersifat Cauchy dan dengan demikian konvergen.  Taksiran dalam
proposisi diperoleh dengan mengambil $m$ menuju tak hingga dalam
\eqref{impropriemann:eqseries} dengan $\ell = k$.

Sebaliknya, andaikan $\int_k^\infty f$ divergen.
Karena $f$ tak negatif, berdasarkan \propref{impropriemann:possimp},
barisan $\{ \int_k^m f \}_{m=k}^\infty$ divergen menuju tak hingga.
Dengan menggunakan
\eqref{impropriemann:eqseries} dengan $\ell = k$, kita memperoleh
\begin{equation*}
\int_k^m f
\leq
\sum_{n=k}^{m-1} f(n) .
\end{equation*}
Karena ruas kiri menuju tak hingga ketika $m \to \infty$, demikian pula
ruas kanan.
\end{proof}

\begin{example}
Uji integral tidak hanya dapat digunakan untuk menunjukkan bahwa suatu deret konvergen, tetapi juga untuk
menaksir jumlahnya dengan ketelitian sebarang.
Mari kita tunjukkan bahwa $\sum_{n=1}^\infty \frac{1}{n^2}$ konvergen.  Setelah itu, kita
taksir jumlahnya dengan galat kurang dari 0.01.  Karena deret ini merupakan deret-$p$ untuk
$p=2$, kita telah membuktikan bahwa deret tersebut konvergen (anggap saja kita belum mengetahuinya),
tetapi kita baru menaksir jumlahnya secara kasar.

Teorema dasar kalkulus menyatakan bahwa untuk semua $k \in \N$,
\begin{equation*}
\int_{k}^\infty \frac{1}{x^2}\,dx = \frac{1}{k} .
\end{equation*}
Khususnya, deret tersebut pasti konvergen.  Namun, kita juga memiliki
\begin{equation*}
\frac{1}{k} = \int_k^\infty \frac{1}{x^2}\,dx
\leq
\sum_{n=k}^\infty \frac{1}{n^2}
\leq
\frac{1}{k^2}
+
\int_k^\infty \frac{1}{x^2}\,dx
=
\frac{1}{k^2}
+
\frac{1}{k} .
\end{equation*}
Dengan menambahkan jumlah parsial hingga $k-1$, kita memperoleh
\begin{equation*}
\frac{1}{k} + \sum_{n=1}^{k-1} \frac{1}{n^2}
\leq
\sum_{n=1}^\infty \frac{1}{n^2}
\leq
\frac{1}{k^2}
+
\frac{1}{k} + \sum_{n=1}^{k-1} \frac{1}{n^2} .
\end{equation*}
Dengan kata lain,
$\nicefrac{1}{k} + \sum_{n=1}^{k-1} \nicefrac{1}{n^2}$ merupakan taksiran
jumlah tersebut dengan galat kurang dari $\nicefrac{1}{k^2}$.  Oleh karena itu, jika kita ingin
menentukan jumlahnya dengan galat kurang dari 0.01, perhatikan bahwa $\nicefrac{1}{{10}^2} = 0.01$.  Kita
memperoleh
\begin{equation*}
1.6397\ldots
\approx
\frac{1}{10} + \sum_{n=1}^{9} \frac{1}{n^2}
\leq
\sum_{n=1}^\infty \frac{1}{n^2}
\leq
\frac{1}{100}
+
\frac{1}{10} + \sum_{n=1}^{9} \frac{1}{n^2}
\approx
1.6497\ldots .
\end{equation*}
Jumlah sebenarnya adalah $\nicefrac{\pi^2}{6} \approx 1.6449\ldots$.
\end{example}

\subsection{Latihan}

\begin{exercise}
Selesaikan pembuktian \propref{impropriemann:ptest}.
\end{exercise}

\begin{exercise}
Tentukan semua $a \in \R$ yang membuat $\sum_{n=1}^\infty e^{an}$ konvergen.
Ketika deret tersebut konvergen, tentukan batas atas bagi jumlahnya.
\end{exercise}

\begin{exercise}
\leavevmode
\begin{enumerate}[a)]
\item
Taksir $\sum_{n=1}^\infty \frac{1}{n(n+1)}$ dengan galat kurang dari 0.01
menggunakan uji integral.
\item
Hitung jumlah deret tersebut secara eksak
dan bandingkan.  Petunjuk: Deretnya bersifat teleskopik.
\end{enumerate}
\end{exercise}

\begin{exercise}
Buktikan bahwa
\begin{equation*}
\int_{-\infty}^\infty \babs{\operatorname{sinc}(x)}\,dx = \infty .
\end{equation*}
Petunjuk: Sekali lagi, cukup tunjukkan hal ini pada satu sisi saja.
\end{exercise}

\begin{exercise}
Dapatkah Anda menafsirkan
\begin{equation*}
\int_{-1}^1 \frac{1}{\sqrt{\sabs{x}}}\,dx
\end{equation*}
sebagai integral tak wajar?  Jika ya, hitung nilainya.
\end{exercise}

\begin{exercise}
Ambil $f \colon [0,\infty) \to \R$ yang terintegralkan secara Riemann pada
setiap interval $[0,b]$ dan sedemikian sehingga terdapat $M$, $a$, dan $T$
dengan $\babs{f(t)} \leq M e^{at}$ untuk semua $t \geq T$.  Tunjukkan bahwa
\emph{\myindex{transformasi Laplace}} dari $f$ ada.  Artinya, untuk
setiap $s > a$, integral berikut konvergen:
\begin{equation*}
F(s) \coloneqq \int_{0}^\infty f(t) e^{-st} \,dt .
\end{equation*}
\end{exercise}

\begin{exercise}
Misalkan $f \colon \R \to \R$ merupakan fungsi yang terintegralkan secara Riemann
pada setiap interval $[a,b]$ dan sedemikian
sehingga $\int_{-\infty}^\infty \babs{f(x)}\,dx < \infty$.  Tunjukkan bahwa
\emph{\myindex{transformasi sinus dan kosinus Fourier}}
ada.  Artinya, untuk setiap $\omega \geq 0$,
integral-integral berikut konvergen:
\begin{equation*}
F^s(\omega) \coloneqq \frac{1}{\pi} \int_{-\infty}^\infty f(t) \sin(\omega t) \,dt ,
\qquad
F^c(\omega) \coloneqq \frac{1}{\pi} \int_{-\infty}^\infty f(t) \cos(\omega t) \,dt .
\end{equation*}
Selain itu, tunjukkan bahwa $F^s$ dan $F^c$ merupakan fungsi terbatas.
\end{exercise}

\begin{exercise}
Andaikan $f \colon [0,\infty) \to \R$ terintegralkan secara Riemann pada setiap interval
$[0,b]$.  Tunjukkan bahwa $\int_0^\infty f$ konvergen jika dan hanya jika
untuk setiap $\epsilon > 0$ terdapat $M$ sedemikian sehingga jika $M \leq a < b$,
maka $\babs{\int_a^b f} < \epsilon$.
\end{exercise}

\begin{exercise}
Andaikan $f \colon [0,\infty) \to \R$ tak negatif dan
\emph{menurun}.  Buktikan:
\begin{enumerate}[a)]
\item
Jika $\int_0^\infty f < \infty$, maka $\lim\limits_{x\to\infty} f(x) = 0$.
\item
Konversnya tidak berlaku.
\end{enumerate}
\end{exercise}

\begin{exercise}
Temukan contoh fungsi kontinu \emph{tak terbatas} $f \colon
[0,\infty) \to \R$ yang tak negatif dan sedemikian sehingga $\int_0^\infty f < \infty$.
Perhatikan bahwa $\lim_{x\to\infty} f(x)$ tidak akan ada; bandingkan dengan
latihan sebelumnya.
Petunjuk: Pada setiap interval $[k,k+1]$, $k \in \N$, definisikan fungsi yang
integralnya pada interval tersebut kurang dari, misalnya, $2^{-k}$.
\end{exercise}

\begin{exercise}[Lebih menantang]
Temukan contoh fungsi $f \colon [0,\infty) \to \R$ yang terintegralkan pada semua
interval sedemikian sehingga $\lim_{n\to\infty} \int_0^n f$ ada sebagai
limit barisan (jadi $n \in \N$), tetapi
$\int_0^\infty f$ tidak ada.
Petunjuk: Untuk semua $n\in \N$, bagi $[n,n+1]$ menjadi dua bagian sama panjang.  Pada satu bagian,
buat fungsi tersebut negatif; pada bagian lainnya, buat fungsi tersebut positif.
\end{exercise}

\begin{exercise}
Andaikan $f \colon [1,\infty) \to \R$ sedemikian sehingga
$g(x) \coloneqq x^2 f(x)$ merupakan fungsi terbatas. Buktikan bahwa
$\int_1^\infty f$ konvergen.
\end{exercise}

\begin{exnote}
Terkadang kita ingin memberikan nilai kepada integral yang biasanya
bahkan tidak dapat ditafsirkan sebagai integral tak wajar,
misalnya $\int_{-1}^1 \nicefrac{1}{x}\,dx$.
Andaikan $f \colon [a,b] \to \R$ merupakan fungsi dan $a < c < b$,
dengan $f$ terintegralkan secara Riemann pada interval
$[a,c-\epsilon]$ dan $[c+\epsilon,b]$ untuk semua
$0 < \epsilon < \min\{c-a,b-c\}$.
Definisikan
\emph{\myindex{nilai utama Cauchy}} dari $\int_a^b f$ sebagai
\begin{equation*}
p.v.\!\int_a^b f \coloneqq \lim_{\epsilon\to 0^+}
\left(
\int_a^{c-\epsilon} f + 
\int_{c+\epsilon}^b f
\right) ,
\end{equation*}
jika limit tersebut ada.
\end{exnote}

%\begin{samepage}
\begin{exercise}
\leavevmode
\begin{enumerate}[a)]
\item
Hitung $p.v.\!\int_{-1}^1 \nicefrac{1}{x}\,dx$.
\item
Hitung
$\lim_{\epsilon\to 0^+}
( \int_{-1}^{-\epsilon} \nicefrac{1}{x}\,dx +
\int_{2\epsilon}^1 \nicefrac{1}{x}\,dx )$ dan tunjukkan bahwa hasilnya tidak sama
dengan nilai utama.
\item
Tunjukkan bahwa jika $f$ terintegralkan pada $[a,b]$, maka
$p.v.\!\int_a^b f = \int_a^b f$ (untuk sebarang $c \in (a,b)$).
\item
Andaikan $f \colon [-1,1] \to \R$
merupakan fungsi ganjil ($f(-x)=-f(x)$) yang terintegralkan pada
$[-1,-\epsilon]$ dan $[\epsilon,1]$ untuk semua $0 < \epsilon < 1$.
Buktikan bahwa
$p.v.\!\int_{-1}^1 f = 0$.
\item
Andaikan
$f \colon [-1,1] \to \R$ kontinu dan diferensiabel di 0.  Tunjukkan bahwa
$p.v.\!\int_{-1}^1 \frac{f(x)}{x}\,dx$ ada.
\end{enumerate}
\end{exercise}
%\end{samepage}

\begin{samepage}
\begin{exercise}
Misalkan $f \colon \R \to \R$ dan
$g \colon \R \to \R$ merupakan fungsi kontinu, dengan
$g(x) = 0$ untuk semua $x \notin [a,b]$ bagi suatu interval $[a,b]$.
\begin{enumerate}[a)]
\item
Tunjukkan bahwa
\emph{\myindex{konvolusi}}
\begin{equation*}
(g * f)(x) \coloneqq \int_{-\infty}^\infty f(t)g(x-t)\,dt 
\end{equation*}
terdefinisi dengan baik untuk semua $x \in \R$.
\item
Andaikan $\int_{-\infty}^\infty \babs{f(x)}\,dx < \infty$.  Buktikan bahwa
\begin{equation*}
\lim_{x \to -\infty} (g * f)(x) = 0, \qquad \text{dan} \qquad
\lim_{x \to \infty} (g * f)(x) = 0 .
\end{equation*}
\end{enumerate}
\end{exercise}
\end{samepage}
